% 06_results.tex -- the rehauled paper's section, as bullets to expand into prose (2026-10-15).
% Sources and pointers behind every number: writeup/PAPER_REHAUL_OUTLINE_2026-10-15.md, section 6.
% House rules: every number in the final prose comes from a generated macro (writeup/generated/*.tex), never
% typed; park (comment out with a dated note) rather than delete.
% Macros to use here: close_main_numbers.tex (\cmMaster*, \cmTrVsRidge*, \cmLd*), fig_qlike_sharpe_numbers.tex (\fqs*); the paired table needs a new pairs CSV and macros

% 2026-10-15 formula pass: each formula checked against the code; disagreements are noted in comments.
% 2026-10-15 (notation pass): symbols and terms per writeup/NOTATION_GUIDE_2026-10-15.md (incl. section 6b). Old -> new:
% RV_t (winsorized) -> \RV_t = \RV_{t,c}, the actual RV; B_t -> s^2_{t,c}; N -> T_oos, sums over \mathcal T_oos;
% lag l -> q; DM -> \DM; (a)/(ref) -> A/B; all / base -> Full / HAR, the forecast f_{m,\mathcal X,\mathcal I};
% per-bar / pooled -> closing-interval / pooled intraday model; HAR + calendar -> HAR benchmark; family -> model
% class; reference OLS -> HAR benchmark; Sharpe -> \SR; 1{.} -> \mathbf{1}{.}. Typed numbers that have an
% exact generated macro now use it. Still typed, no macro yet: the per-model QLIKE levels, % changes and DM values
% (lasso 0.0914 / -3.4 / -1.22, ...; they are in tab:rh_axis1, tab:rh_axis2, tab:rh_rank), the 83--91 % squared-error
% share, and the pooled ranges. All of them are pre-switch values (scored against the winsorized target).
% 2026-10-15 (consolidation): NW lag q -> q_NW (q_t is the position); the OLS on the HAR benchmark is "the
% benchmark forecast" (the HAR benchmark is the predictor set); every number re-scored on the actual RV now comes
% from \rhr... macros (scoring agent #10; the \rhrQ / \rhrVsBm / \rhrAf / \rhrPoolVsClose families). No typed result
% remains in this section. Sign convention for the estimation-sample comparison: pooled intraday minus
% closing-interval, in Table T2 and (since this pass) in Table E3.
% 2026-10-15 (user, restructure): section renamed "Experiments and results" (label sec:rh_results kept); it opens
% with "Experimental setup" (sec:rh_exp_setup): samples and windows (sec:rh_exp_samples; values moved from Methods'
% Problem formulation), training schemes in practice (sec:rh_exp_schemes), tuned hyperparameters
% (sec:rh_methods_hyper) and tuning procedures (sec:rh_methods_tuning), the last two moved from Methods with their
% labels; their tables (tab:rh_dm_tuning, tab:rh_dm_hyper, tab:rh_dm_hyper_search) moved to 06_results_floats.tex.
% Phase 2: typed 255, 24{,}000, 125, "two" -> \rhdmRfTrial, \rhdmPooledWindow, \rhdmValTail, \rhdmRetuneFolds.
% 2026-10-15 (finish): the DM horizon symbol h dropped ("one-step-ahead forecasts"; h = 1 substituted in q_NW); the
% sign strategy described at its first use (the Introduction is parked); MinGRU bullets and the MinGRU entry of the
% accuracy ranking parked (user: NN content parked for now).
% 2026-10-15 (user, restructure): "I need the last 30 min scheme to be moved to the end of the economic application
% section as a final cherry on top instead of competing with the exogenous feature signal story". This section now
% presents the pooled intraday models only (the paper's training scheme); every closing-interval (last-30-minutes)
% result is parked in place and restated in the final subsection of Economic value (sec:rh_econ_last30). The old
% subsection "Closing-interval versus pooled intraday models" is parked whole (its label moved there).
% 2026-10-15 (consolidation): parked -- the number files are loaded once in the preamble of paper_rehaul.tex.
% \ifdefined\rhrDays\else% AUTO-GENERATED by writeup/make_rehaul_results.py -- do not edit.
% Numbers for the captions of writeup/rehaul/06_results_floats.tex (plain f^2 B, 866 days, DM lag 6 HLN).
\providecommand{\rhrDays}{866}  % trade days scored
\providecommand{\rhrNNPooledSeeds}{3}  % seeds averaged in the pooled MinGRU forecast
\providecommand{\rhrOlsQ}{0.0946}  % QLIKE of the per-bar OLS on HAR + calendar
\providecommand{\rhrAfPairs}{51}  % model-set input-set pairs (all features vs HAR + cal.)
\providecommand{\rhrAfLower}{38}  % ... QLIKE lower with all features
\providecommand{\rhrAfLowerSig}{2}  % ... and DM p < 0.05
\providecommand{\rhrAfMseHigher}{51}  % ... variance-space MSE higher with all features
\providecommand{\rhrAfMseHigherSig}{0}  % ... and DM (squared error) p < 0.05
\providecommand{\rhrNModelSet}{112}  % forecasts in the master table's model set
\providecommand{\rhrRankN}{120}  % forecasts ranked in the full table
\providecommand{\rhrNSet}{111}  % model-set forecasts other than the OLS
\providecommand{\rhrNLower}{7}  % ... with lower QLIKE than the OLS
\providecommand{\rhrNLowerSig}{0}  % ... and DM p < 0.05
\providecommand{\rhrNWorseSig}{48}  % ... significantly worse than the OLS (DM p < 0.05)
\providecommand{\rhrLinBaseMaxPct}{0.03}  % largest |% change| of the per-bar ridge / lasso / enet on HAR + cal. vs the OLS
\providecommand{\rhrSETopFiveOls}{88}  % % of the OLS's total squared error on its 5 largest days
\providecommand{\rhrSETopFiveMin}{85}  % min over the four F5 forecasts: % of squared error on the 5 largest days
\providecommand{\rhrSETopFiveMax}{92}  % max over the four F5 forecasts
\providecommand{\rhrQLTopFiveMin}{7}  % min over the four F5 forecasts: % of total QLIKE on the 5 largest days
\providecommand{\rhrQLTopFiveMax}{9}  % max over the four F5 forecasts
\providecommand{\rhrCumLassoTotal}{\ensuremath{-}2.76}  % lasso, all features: sum over days of the QLIKE difference vs the OLS
\providecommand{\rhrCumLassoWindow}{\ensuremath{+}1.60}  % ... of which in 24 Feb - 30 Apr 2020
\providecommand{\rhrCumXgbTotal}{\ensuremath{-}0.98}  % XGBoost (daily refit), all features: sum over days
\providecommand{\rhrCumXgbWindow}{\ensuremath{+}0.30}  % ... of which in 24 Feb - 30 Apr 2020
\providecommand{\rhrCumNnTotal}{\ensuremath{+}0.16}  % MinGRU, all features: sum over days
\providecommand{\rhrCumNnWindow}{\ensuremath{+}0.78}  % ... of which in 24 Feb - 30 Apr 2020
\fi
% \ifdefined\fqsPairs\else% AUTO-GENERATED by writeup/make_fig_qlike_sharpe.py -- do not edit.
% Figure writeup/figures/fig_qlike_sharpe.pdf: QLIKE against Sharpe (mid), HAR + calendar vs all
% features, every model of results/close_master_table/paper_model_set.csv on both input sets
% (per-bar linear = the pre-tuned subpt_* rows; LSTMs excluded); numbers from master_table.csv
% and master_table_daily.parquet.  Changes are all features minus HAR + calendar, on the same days.
% DM: dm_test on the daily QLIKE difference; t: hac_inference.llsw_test on the daily mid-return
% difference (614 post-warmup days, bandwidth 33, fixed-b p); Sharpe changes: Ledoit-Wolf p;
% ES(c = 1): close_murphy_splits_headline.elementary on f / iv, y / iv.  Point estimates.
% pooled pairs: family pooled twin
\newcommand{\fqsPairs}{42}  % models on both input sets (pairs drawn)
\newcommand{\fqsPairsPerBarLinear}{3}  % of which per-bar linear (pre-tuned)
\newcommand{\fqsPairsPooled}{3}  % of which pooled 48-bar linear
\newcommand{\fqsPairsTree}{36}  % of which per-bar trees
\newcommand{\fqsPoints}{85}  % points drawn: both ends of every pair + the HAR + cal. OLS
\newcommand{\fqsDays}{866}  % trade days
\newcommand{\fqsQLower}{31}  % QLIKE lower with all features
\newcommand{\fqsQLowerSig}{3}  % ... and DM significant at 5 %
\newcommand{\fqsQHigher}{11}  % QLIKE higher with all features
\newcommand{\fqsQHigherSig}{0}  % ... and DM significant at 5 %
\newcommand{\fqsQLowerSigLagFourteen}{2}  % sensitivity: QLIKE lower, DM at lag 14 significant
\newcommand{\fqsQHigherSigLagFourteen}{0}  % sensitivity: QLIKE higher, DM at lag 14 significant
\newcommand{\fqsShHigher}{41}  % Sharpe (mid) higher with all features
\newcommand{\fqsShHigherSig}{6}  % ... and Ledoit-Wolf p < 0.05 (614 post-warmup days)
\newcommand{\fqsShLower}{1}  % Sharpe (mid) lower with all features
\newcommand{\fqsShLowerSig}{0}  % ... and Ledoit-Wolf p < 0.05 (614 post-warmup days)
\newcommand{\fqsQLowerShHigher}{30}  % both: QLIKE lower and Sharpe higher
\newcommand{\fqsMedDQ}{\ensuremath{-}0.0016}  % median QLIKE change (all features minus HAR + cal.)
\newcommand{\fqsMedDQPct}{\ensuremath{-}1.6}  % median QLIKE change, % of the HAR + cal. QLIKE
\newcommand{\fqsMedDSh}{\ensuremath{+}0.53}  % median Sharpe (mid) change
\newcommand{\fqsESLower}{29}  % elementary score at the threshold (c = 1) lower with all features
\newcommand{\fqsCorrES}{0.07}  % Pearson corr. across pairs: ES(c = 1) change vs HAC t of the P&L change
\newcommand{\fqsCorrQ}{0.12}  % Pearson corr. across pairs: QLIKE change vs HAC t of the P&L change
\newcommand{\fqsInShade}{4}  % all-features points in the shaded quadrant
\newcommand{\fqsInShadeBase}{3}  % HAR + cal. points in the shaded quadrant
\newcommand{\fqsOlsQ}{0.0967}  % QLIKE of the HAR + cal. OLS
\newcommand{\fqsOlsSh}{1.06}  % Sharpe (mid) of the HAR + cal. OLS
\newcommand{\fqsDMLag}{6}  % Newey-West lag of the DM statistic (the master table's)
\newcommand{\fqsHACS}{33}  % bandwidth S of the t of the mean P&L change (LLSW 2018, 614 post-warmup days)
\newcommand{\fqsSigPct}{5}  % two-sided level of the counts, %
\newcommand{\fqsQuadBetterBoth}{30}  % change figure: lower QLIKE, higher Sharpe
\newcommand{\fqsQuadShOnly}{11}  % change figure: higher QLIKE, higher Sharpe
\newcommand{\fqsQuadQOnly}{1}  % change figure: lower QLIKE, lower Sharpe
\newcommand{\fqsQuadWorseBoth}{0}  % change figure: higher QLIKE, lower Sharpe
\fi
% \ifdefined\cmNdays\else% AUTO-GENERATED by writeup/make_table_close_main.py -- do not edit.
% Numbers quoted in the prose of writeup/sections/results_close_option.tex; each is read
% from results/close_master_table/master_table.csv and the sources in the script.
% 2026-10-06: recalibration removed / intervals removed, user decision -- the plain forecast
% f^2 B, point estimates; ...T = HAC t of the daily mid-return difference, ...DM = Diebold-
% Mariano statistic; a ...Above / ...Below count counts point estimates, its ...Sig twin those
% significant at 5 % (DM for QLIKE, HAC t for Sharpe).
\newcommand{\cmNdays}{866}
\newcommand{\cmNrows}{22}
\newcommand{\cmShortMean}{0.014}
\newcommand{\cmShortShMid}{0.20}
\newcommand{\cmShortShX}{\ensuremath{-}0.27}
\newcommand{\cmShortT}{\ensuremath{-}1.71}
\newcommand{\cmPaperShMin}{0.43}
\newcommand{\cmPaperShMax}{1.50}
\newcommand{\cmPaperXMin}{\ensuremath{-}0.05}
\newcommand{\cmPaperXMax}{1.01}
\newcommand{\cmPaperVsShortMin}{+0.22}
\newcommand{\cmPaperVsShortMax}{+1.30}
\newcommand{\cmPaperVsShortAboveN}{8}
\newcommand{\cmPaperVsShortAboveSigN}{three}
\newcommand{\cmPaperVsShortAboveSigWho}{LightGBM, XGBoost, lasso (causally tuned)}
\newcommand{\cmPaperSameMin}{82}
\newcommand{\cmSellPaperMin}{72}
\newcommand{\cmSellPaperMax}{83}
\newcommand{\cmPaperSameMax}{96}
\newcommand{\cmRefQ}{0.1092}
\newcommand{\cmRefShMid}{1.28}
\newcommand{\cmRefShX}{0.80}
\newcommand{\cmRefBuy}{27}
\newcommand{\cmRefVsShort}{+1.08}
\newcommand{\cmRefVsShortT}{+1.71}
\newcommand{\cmFomcSh}{1.36}
\newcommand{\cmFomcD}{+0.08}
\newcommand{\cmFomcDT}{+0.44}
\newcommand{\cmLassoFShMid}{1.06}
\newcommand{\cmLassoFVsRef}{\ensuremath{-}0.22}
\newcommand{\cmLassoFVsRefT}{\ensuremath{-}0.75}
\newcommand{\cmLassoFSame}{96}
\newcommand{\cmPaperFarLabel}{baseline (HAR + calendar OLS)}
\newcommand{\cmPaperFarVsRef}{\ensuremath{-}0.85}
\newcommand{\cmPaperFarVsRefT}{\ensuremath{-}1.49}
\newcommand{\cmPaperBestLabel}{lasso (causally tuned)}
\newcommand{\cmPaperBestSh}{1.50}
\newcommand{\cmNlFamN}{five}
\newcommand{\cmNlMasterN}{90}
\newcommand{\cmPoolShMin}{0.91}
\newcommand{\cmPoolShMax}{1.44}
\newcommand{\cmSubShMin}{1.37}
\newcommand{\cmSubShMax}{1.92}
\newcommand{\cmNlShMin}{1.23}
\newcommand{\cmNlShMax}{1.81}
\newcommand{\cmTabVsRefN}{21}
\newcommand{\cmTabVsRefAbove}{15}
\newcommand{\cmTabVsRefBelow}{6}
\newcommand{\cmTabVsRefAboveSig}{0}
\newcommand{\cmTabVsRefBelowSig}{0}
\newcommand{\cmMasterN}{112}
\newcommand{\cmMasterNPaper}{9}
\newcommand{\cmMasterNBase}{52}
\newcommand{\cmMasterMedQBase}{0.1033}
\newcommand{\cmMasterMedShBase}{0.73}
\newcommand{\cmMasterMedShXBase}{0.25}
\newcommand{\cmMasterNAll}{51}
\newcommand{\cmMasterMedQAll}{0.1014}
\newcommand{\cmMasterMedShAll}{1.27}
\newcommand{\cmMasterMedShXAll}{0.79}
\newcommand{\cmMasterSetsN}{two}
\newcommand{\cmMasterVsRefN}{111}
\newcommand{\cmMasterVsRefAbove}{32}
\newcommand{\cmMasterVsRefBelow}{79}
\newcommand{\cmMasterVsRefAboveSig}{0}
\newcommand{\cmMasterVsRefBelowSig}{0}
\newcommand{\cmMasterVsShortAbove}{112}
\newcommand{\cmMasterVsShortBelow}{0}
\newcommand{\cmMasterVsShortAboveSig}{20}
\newcommand{\cmMasterVsShortBelowSig}{0}
\newcommand{\cmMasterVsShortN}{112}
\newcommand{\cmMasterQVsRefBelow}{99}
\newcommand{\cmMasterQVsRefAbove}{12}
\newcommand{\cmMasterQVsRefBelowSig}{23}
\newcommand{\cmMasterQVsRefAboveSig}{1}
\newcommand{\cmMasterMedQ}{0.1028}
\newcommand{\cmMasterMedSh}{1.08}
\newcommand{\cmMasterMedShX}{0.59}
\newcommand{\cmMasterShMin}{0.29}
\newcommand{\cmMasterShMax}{1.92}
\newcommand{\cmMasterXMin}{\ensuremath{-}0.19}
\newcommand{\cmMasterXMax}{1.43}
\newcommand{\cmMasterVsShortMed}{+0.87}
\newcommand{\cmMasterRankCorr}{\ensuremath{-}0.20}
\newcommand{\cmMasterRankCorrX}{\ensuremath{-}0.20}
\newcommand{\cmSellMasterMin}{66}
\newcommand{\cmSellMasterMax}{83}
\newcommand{\cmSellMasterMed}{70}
\newcommand{\cmTopShLabel}{per-bar elastic net}
\newcommand{\cmTopShInputs}{\texttt{all\_features}}
\newcommand{\cmTopShMid}{1.92}
\newcommand{\cmXLossMin}{0.48}
\newcommand{\cmXLossMax}{0.50}
\newcommand{\cmLdSpec}{+0.16}
\newcommand{\cmLdSpecQ}{\ensuremath{-}4.3}
\newcommand{\cmLdSpecT}{+0.34}
\newcommand{\cmLdSpecDM}{\ensuremath{-}1.73}
\newcommand{\cmLdPerBarBase}{+0.46}
\newcommand{\cmLdPerBarBaseQ}{\ensuremath{-}12.0}
\newcommand{\cmLdPerBarBaseT}{+1.05}
\newcommand{\cmLdPerBarBaseDM}{\ensuremath{-}3.37}
\newcommand{\cmLdPerBarAll}{\ensuremath{-}0.03}
\newcommand{\cmLdPerBarAllQ}{\ensuremath{-}6.7}
\newcommand{\cmLdPerBarAllT}{\ensuremath{-}0.07}
\newcommand{\cmLdPerBarAllDM}{\ensuremath{-}1.61}
\newcommand{\cmLdIn}{+0.05}
\newcommand{\cmLdInQ}{+0.8}
\newcommand{\cmLdInT}{+0.09}
\newcommand{\cmLdInDM}{+0.18}
\newcommand{\cmLdEstMin}{\ensuremath{-}0.15}
\newcommand{\cmLdEstMax}{+0.50}
\newcommand{\cmLdEstMinLabel}{\texttt{subtree\_rf\_all\_features}}
\newcommand{\cmLdEstMaxLabel}{\texttt{subpt\_enet\_all\_features}}
\newcommand{\cmLdEstMinWho}{random forest}
\newcommand{\cmLdEstMaxWho}{elastic net}
\newcommand{\cmLdN}{21}
\newcommand{\cmLdNPath}{12}
\newcommand{\cmLdNCampaign}{9}
\newcommand{\cmLdResolved}{none}
\newcommand{\cmLdQResolved}{three}
\newcommand{\cmLdCDedup}{+0.00}
\newcommand{\cmLdCDedupQ}{+0.0}
\newcommand{\cmLdCMaskTten}{\ensuremath{-}0.21}
\newcommand{\cmLdCMaskTtenQ}{\ensuremath{-}1.0}
\newcommand{\cmLdCMaskTtenT}{\ensuremath{-}0.69}
\newcommand{\cmLdCMaskTtenDM}{\ensuremath{-}1.04}
\newcommand{\cmLdCMaskTone}{\ensuremath{-}0.23}
\newcommand{\cmLdCMaskToneQ}{+1.4}
\newcommand{\cmLdCMaskToneT}{\ensuremath{-}0.55}
\newcommand{\cmLdCMaskToneDM}{+1.36}
\newcommand{\cmLdCMaskRsten}{\ensuremath{-}0.08}
\newcommand{\cmLdCMaskRstenQ}{+2.8}
\newcommand{\cmLdCMaskRstenT}{\ensuremath{-}0.32}
\newcommand{\cmLdCMaskRstenDM}{+1.42}
\newcommand{\cmLdCMaskRsone}{\ensuremath{-}0.05}
\newcommand{\cmLdCMaskRsoneQ}{+0.2}
\newcommand{\cmLdCMaskRsoneT}{\ensuremath{-}0.17}
\newcommand{\cmLdCMaskRsoneDM}{+0.12}
\newcommand{\cmLdCCadLgbm}{\ensuremath{-}0.04}
\newcommand{\cmLdCCadLgbmQ}{\ensuremath{-}2.2}
\newcommand{\cmLdCCadLgbmT}{\ensuremath{-}0.08}
\newcommand{\cmLdCCadLgbmDM}{\ensuremath{-}1.12}
\newcommand{\cmLdCTuneRs}{+0.28}
\newcommand{\cmLdCTuneRsQ}{+3.0}
\newcommand{\cmLdCTuneRsT}{+0.58}
\newcommand{\cmLdCTuneRsDM}{+1.12}
\newcommand{\cmLdCTuneOp}{\ensuremath{-}0.10}
\newcommand{\cmLdCTuneOpQ}{+7.3}
\newcommand{\cmLdCTuneOpT}{\ensuremath{-}0.31}
\newcommand{\cmLdCTuneOpDM}{+2.83}
\newcommand{\cmLdCTuneOpPer}{\ensuremath{-}0.32}
\newcommand{\cmLdCTuneOpPerQ}{+2.5}
\newcommand{\cmLdCTuneOpPerT}{\ensuremath{-}0.65}
\newcommand{\cmLdCTuneOpPerDM}{+0.75}
\newcommand{\cmPNewBase}{22}
\newcommand{\cmPOldBase}{34}
\newcommand{\cmPNewAll}{628}
\newcommand{\cmPOldAll}{640}
\newcommand{\cmNDup}{12}
\newcommand{\cmDedupN}{37}
\newcommand{\cmDedupQMax}{0.05}
\newcommand{\cmDedupPosMax}{one}
\newcommand{\cmKeptBase}{17}
\newcommand{\cmKeptBaseMin}{17}
\newcommand{\cmKeptBaseMax}{17}
\newcommand{\cmKeptAll}{369}
\newcommand{\cmKeptAllMin}{360}
\newcommand{\cmKeptAllMax}{378}
\newcommand{\cmMaskN}{16}
\newcommand{\cmMaskQMed}{0.4}
\newcommand{\cmMaskQAbove}{10}
\newcommand{\cmMaskShBelow}{14}
\newcommand{\cmMaskQAboveSig}{none}
\newcommand{\cmMaskShBelowSig}{one}
\newcommand{\cmTrVsRidgeN}{48}
\newcommand{\cmTrSetsN}{two}
\newcommand{\cmMaskSetsN}{two}
\newcommand{\cmTrVsRidgeQAbove}{46}
\newcommand{\cmTrVsRidgeShBelow}{41}
\newcommand{\cmTrVsRidgeQAboveSig}{21}
\newcommand{\cmTrVsRidgeShBelowSig}{one}
\newcommand{\cmCadN}{six}
\newcommand{\cmCadQBelow}{five}
\newcommand{\cmCadQBelowSig}{none}
\newcommand{\cmCadShUp}{\ensuremath{-}0.43}
\newcommand{\cmCadShUpT}{\ensuremath{-}1.03}
\newcommand{\cmRsQAbove}{6}
\newcommand{\cmRsQAboveSig}{two}
\newcommand{\cmRsN}{6}
\newcommand{\cmOpN}{24}
\newcommand{\cmOpQAbove}{24}
\newcommand{\cmOpQAboveSig}{13}
\newcommand{\cmTpN}{18}
\newcommand{\cmTpQBelow}{12}
\newcommand{\cmTpQBelowSig}{one}
\newcommand{\cmBudN}{24}
\newcommand{\cmBudQAbove}{17}
\newcommand{\cmBudQAboveSig}{two}
\newcommand{\cmBudQBelow}{7}
\newcommand{\cmBudLastMin}{27}
\newcommand{\cmBudLastMax}{37}
\newcommand{\cmBudLastExch}{20}
\newcommand{\cmNlBestSh}{1.81}
\newcommand{\cmNlBestLabel}{tuned per-bar LightGBM}
\newcommand{\cmNlBestInputs}{\texttt{all\_features}}
\newcommand{\cmNlBestVsRef}{+0.53}
\newcommand{\cmNlBestVsRefT}{+1.04}
\newcommand{\cmNlVsRefAbove}{19}
\newcommand{\cmNlVsRefBelow}{71}
\newcommand{\cmNlVsRefAboveSig}{0}
\newcommand{\cmNlVsRefBelowSig}{0}
\newcommand{\cmClassChunks}{505}
\newcommand{\cmClassPts}{20}
\newcommand{\cmClassLgbmPick}{18}
\newcommand{\cmClassMaxRelMin}{4.5}
\newcommand{\cmClassMaxRelMax}{5.5}
\newcommand{\cmSvN}{31}
\newcommand{\cmSvSetN}{112}
\newcommand{\cmSvSetPairsN}{3472}
\newcommand{\cmSvSetAboveMid}{755}
\newcommand{\cmSvSetAboveMidSig}{29}
\newcommand{\cmSvSetStraddleFirstMid}{4}
\newcommand{\cmSvSetWingN}{672}
\newcommand{\cmSvSetWingMidBelow}{668}
\newcommand{\cmSvSetAboveX}{410}
\newcommand{\cmSvSetAboveXSig}{0}
\newcommand{\cmSvSetStraddleFirstX}{3}
\newcommand{\cmSvSetWingXBelow}{672}
\newcommand{\cmSvSetBestMidLabel}{call vertical, w50}
\newcommand{\cmSvSetBestMidMed}{+0.37}
\newcommand{\cmSvSetBestMidAboveN}{94}
\newcommand{\cmSvSetBestMidAboveSig}{0}
\newcommand{\cmSvCostStraddle}{2.9}
\newcommand{\cmSvCostIbfMin}{4.0}
\newcommand{\cmSvCostIbfMax}{6.1}
\newcommand{\cmMasterPosTotalN}{112}
\newcommand{\cmMasterTopTwentyShareMed}{99}
\newcommand{\cmMasterTopTwentyShareMin}{59}
\newcommand{\cmMasterTopTwentyShareMax}{338}
\newcommand{\cmMasterTopTwentyBuysN}{112}
\newcommand{\cmMasterBottomTwentySellsN}{112}
\newcommand{\cmMasterDiffMed}{+53.5}
\newcommand{\cmMasterDiffPosN}{112}
\newcommand{\cmMasterDiffExTenMed}{\ensuremath{-}30.0}
\newcommand{\cmMasterDiffExTenHalfN}{112}
\newcommand{\cmMasterDiffExTenNegN}{102}
\newcommand{\cmBContamMaxDSh}{0.38}
\newcommand{\cmBContamMaxDShX}{0.38}
\newcommand{\cmBContamMedDSh}{0.08}
\newcommand{\cmBContamPaperMaxDSh}{0.21}
\newcommand{\cmBContamPaperMaxDShX}{0.22}
\newcommand{\cmBContamPaperMedDSh}{0.08}
\newcommand{\cmBContamDays}{94}
\newcommand{\cmCensIvMaxDSh}{0.06}
\newcommand{\cmCensIvMaxDShX}{0.08}
\newcommand{\cmCensIvMedDSh}{0.04}
\newcommand{\cmCensIvPaperMaxDSh}{0.05}
\newcommand{\cmCensIvPaperMaxDShX}{0.07}
\newcommand{\cmCensIvPaperMedDSh}{0.04}
\newcommand{\cmCensIvDays}{5}
\newcommand{\cmBContamMaxFlips}{16}
\newcommand{\cmBContamMedFlips}{8}
\newcommand{\cmBContamPaperMaxFlips}{10}
\newcommand{\cmBContamMaxRel}{21}
\newcommand{\cmPtPreSessions}{3947}
\newcommand{\cmPtPreStart}{13~April~2004}
\newcommand{\cmPtValFirst}{6~January~2016}
\newcommand{\cmPtValLast}{2~January~2020}
\newcommand{\cmPtFolds}{four}
\newcommand{\cmPtFoldLen}{250}
\newcommand{\cmPtFitWin}{2{,}000}
\newcommand{\cmPtEmbargo}{25}
\newcommand{\cmPtRatio}{1/100}
\newcommand{\cmPtBlock}{250}
\newcommand{\cmPtRetuneN}{four}
\newcommand{\cmPtRetuneFirst}{23~June~2020}
\newcommand{\cmPtRetuneLast}{16~June~2023}
\newcommand{\cmPtFiredRidgeBase}{one}
\newcommand{\cmPtPenRidgeBase}{\ensuremath{10^{-2}}}
\newcommand{\cmPtFiredRidgeAll}{none}
\newcommand{\cmPtPenRidgeAll}{\ensuremath{10^{3}}}
\newcommand{\cmPtFiredLassoBase}{one}
\newcommand{\cmPtPenLassoBase}{\ensuremath{10^{-6}}}
\newcommand{\cmPtFiredLassoAll}{none}
\newcommand{\cmPtPenLassoAll}{\ensuremath{10^{-2}}}
\newcommand{\cmPtFiredEnetBase}{two}
\newcommand{\cmPtPenEnetBase}{\ensuremath{10^{-6}}}
\newcommand{\cmPtFiredEnetAll}{none}
\newcommand{\cmPtPenEnetAll}{\ensuremath{10^{-2}}}
\newcommand{\cmPtFiredTotal}{4}
\newcommand{\cmPtFiredOf}{24}
\newcommand{\cmPtCausalSwitches}{75}
\newcommand{\cmPtCausalChoices}{165}
\newcommand{\cmPtCausalMarginPass}{3}
\newcommand{\cmPtCausalEdgePicks}{132}
\newcommand{\cmPtCausalPicks}{198}
\newcommand{\cmPtSwitchFlipPct}{8}
\newcommand{\cmPtCausalSwitchesPaper}{24}
\newcommand{\cmPtCausalChoicesPaper}{45}
\newcommand{\cmPtCausalMarginPassPaper}{3}
\newcommand{\cmPtCausalEdgePicksPaper}{34}
\newcommand{\cmPtCausalPicksPaper}{54}
\newcommand{\cmPtSwitchFlipPctPaper}{9}
\newcommand{\cmPoolPtN}{6}
\newcommand{\cmPoolPtQPctMean}{+0.9}
\newcommand{\cmPoolPtQWorse}{5}
\newcommand{\cmPoolPtDMMax}{1.22}
\newcommand{\cmPoolPtTMax}{1.29}
\newcommand{\cmPoolPtSigN}{0}
\newcommand{\cmPoolPtShAllMin}{+0.17}
\newcommand{\cmPoolPtShAllMax}{+0.35}
\newcommand{\cmPoolPtShOtherMin}{\ensuremath{-}0.06}
\newcommand{\cmPoolPtShOtherMax}{\ensuremath{-}0.19}
\newcommand{\cmPoolPtBlocks}{218}
\newcommand{\cmPoolPtCausalSwMin}{137}
\newcommand{\cmPoolPtCausalSwMax}{170}
\newcommand{\cmPoolPtCausalMarginMin}{3}
\newcommand{\cmPoolPtCausalMarginMax}{22}
\newcommand{\cmPoolPtRetunes}{217}
\newcommand{\cmPoolPtRetuneSwMin}{11}
\newcommand{\cmPoolPtRetuneSwMax}{40}
\newcommand{\cmTreeProvRfName}{\texttt{rf\_opt\_00}}
\newcommand{\cmTreeProvRfTrials}{400}
\newcommand{\cmTreeProvRfBudget}{400}
\newcommand{\cmTreeProvRfDevMse}{0.108456}
\newcommand{\cmTreeProvLgbmDevMse}{0.107449}
\newcommand{\cmTreeProvXgbDevMse}{0.108076}
\newcommand{\cmTreeProvXgbTopDevMse}{0.107894}
\newcommand{\cmTreeProvRfTrees}{485}
\newcommand{\cmTreeProvRfDepth}{20}
\newcommand{\cmTreeProvRfFeatures}{0.90}
\newcommand{\cmTreeProvRfSamples}{0.41}
\newcommand{\cmTreeProvRfLeafDev}{12}
\newcommand{\cmTreeProvRfLeaf}{1}
\newcommand{\cmTreeProvLgbmLeafDev}{98}
\newcommand{\cmTreeProvLgbmLeaf}{8}
\newcommand{\cmTreeProvXgbTopName}{\texttt{xgb\_opt\_00}}
\newcommand{\cmTreeProvXgbName}{\texttt{xgb\_opt\_06}}
\newcommand{\cmTreeProvXgbDevRank}{7}
\newcommand{\cmTreeProvXgbArms}{8}
\newcommand{\cmTreeProvXgbPostRank}{8}
\newcommand{\cmTreeProvXgbRecentRank}{8}
\newcommand{\cmTreeProvXgbQRank}{8}
\newcommand{\cmTreeProvXgbPickDate}{14~August~2026}
\newcommand{\cmMasterNLean}{51}
\newcommand{\cmMasterMedQLean}{0.0991}
\newcommand{\cmMasterMedShLean}{1.37}
\newcommand{\cmMasterMedShXLean}{0.88}
\newcommand{\cmLdPerBarLean}{+0.65}
\newcommand{\cmLdPerBarLeanQ}{\ensuremath{-}12.3}
\newcommand{\cmLdPerBarLeanT}{+1.49}
\newcommand{\cmLdPerBarLeanDM}{\ensuremath{-}3.39}
\newcommand{\cmLdInBaseLean}{+0.51}
\newcommand{\cmLdInBaseLeanQ}{\ensuremath{-}3.3}
\newcommand{\cmLdInBaseLeanT}{+1.66}
\newcommand{\cmLdInBaseLeanDM}{\ensuremath{-}1.17}
\newcommand{\cmLdInLeanAll}{\ensuremath{-}0.46}
\newcommand{\cmLdInLeanAllQ}{+4.3}
\newcommand{\cmLdInLeanAllT}{\ensuremath{-}1.15}
\newcommand{\cmLdInLeanAllDM}{+1.50}
%
% 2026-10-06 parked (recalibration removed / intervals removed, user decision):
%   \cmShortShCiMid: interval of always short's Sharpe
%   \cmShortShCiX: interval of always short's Sharpe
%   \cmPaperVsShortAboveWho: the paper's forecasts whose interval vs always short lay above zero
%   \cmPaperVsShortAboveCi: that forecast's interval
%   \cmSellHead: the former headline's sell share (no headline model)
%   \cmHQ: the former headline (no headline model)
%   \cmHQPct: the former headline (no headline model)
%   \cmHQCi: the former headline (no headline model)
%   \cmHDM: the former headline (no headline model)
%   \cmHDMp: the former headline (no headline model)
%   \cmHShMid: the former headline (no headline model)
%   \cmHShX: the former headline (no headline model)
%   \cmHMean: the former headline (no headline model)
%   \cmHHit: the former headline (no headline model)
%   \cmHBuy: the former headline (no headline model)
%   \cmHVsRefMid: the former headline (no headline model)
%   \cmHVsRefX: the former headline (no headline model)
%   \cmHVsShortMid: the former headline (no headline model)
%   \cmHVsShortX: the former headline (no headline model)
%   \cmHVsRefT: the former headline (no headline model)
%   \cmHVsRefP: the former headline (no headline model)
%   \cmHVsShortT: the former headline (no headline model)
%   \cmHVsShortP: the former headline (no headline model)
%   \cmHRankSh: the former headline (no headline model)
%   \cmHRankQ: the former headline (no headline model)
%   \cmVsHN: the former headline (no headline model)
%   \cmVsHAbove: the former headline (no headline model)
%   \cmVsHBelow: the former headline (no headline model)
%   \cmVsHAboveSig: the former headline (no headline model)
%   \cmVsHBelowSig: the former headline (no headline model)
%   \cmVsHLassoSh: the former headline (no headline model)
%   \cmVsHLasso: paired against the former headline (no headline model)
%   \cmVsHLassoT: paired against the former headline (no headline model)
%   \cmVsHLassoSame: paired against the former headline (no headline model)
%   \cmVsHLassoQPct: paired against the former headline (no headline model)
%   \cmVsHLassoDM: paired against the former headline (no headline model)
%   \cmVsHEnet: paired against the former headline (no headline model)
%   \cmVsHEnetT: paired against the former headline (no headline model)
%   \cmVsHEnetSame: paired against the former headline (no headline model)
%   \cmVsHEnetQPct: paired against the former headline (no headline model)
%   \cmVsHEnetDM: paired against the former headline (no headline model)
%   \cmVsHRidgeAll: paired against the former headline (no headline model)
%   \cmVsHRidgeAllT: paired against the former headline (no headline model)
%   \cmVsHRidgeAllSame: paired against the former headline (no headline model)
%   \cmVsHRidgeAllQPct: paired against the former headline (no headline model)
%   \cmVsHRidgeAllDM: paired against the former headline (no headline model)
%   \cmVsHTop: paired against the former headline (no headline model)
%   \cmVsHTopT: paired against the former headline (no headline model)
%   \cmVsHTopSame: paired against the former headline (no headline model)
%   \cmVsHTopQPct: paired against the former headline (no headline model)
%   \cmVsHTopDM: paired against the former headline (no headline model)
%   \cmLdPerBarLive: a step to the former headline (no headline, no live-feasible)
%   \cmLdPerBarLiveQ: a step to the former headline (no headline, no live-feasible)
%   \cmLdPerBarLiveT: a step to the former headline (no headline, no live-feasible)
%   \cmLdPerBarLiveDM: a step to the former headline (no headline, no live-feasible)
%   \cmLdInAll: a step to the former headline (no headline, no live-feasible)
%   \cmLdInAllQ: a step to the former headline (no headline, no live-feasible)
%   \cmLdInAllT: a step to the former headline (no headline, no live-feasible)
%   \cmLdInAllDM: a step to the former headline (no headline, no live-feasible)
%   \cmLdInBase: a step to the former headline (no headline, no live-feasible)
%   \cmLdInBaseQ: a step to the former headline (no headline, no live-feasible)
%   \cmLdInBaseT: a step to the former headline (no headline, no live-feasible)
%   \cmLdInBaseDM: a step to the former headline (no headline, no live-feasible)
%   \cmLdCMaskLstm: an LSTM step of the campaign (LSTMs parked)
%   \cmLdCMaskLstmQ: an LSTM step of the campaign (LSTMs parked)
%   \cmLdCMaskLstmT: an LSTM step of the campaign (LSTMs parked)
%   \cmLdCMaskLstmDM: an LSTM step of the campaign (LSTMs parked)
%   \cmLdCCadLstm: an LSTM step of the campaign (LSTMs parked)
%   \cmLdCCadLstmQ: an LSTM step of the campaign (LSTMs parked)
%   \cmLdCCadLstmT: an LSTM step of the campaign (LSTMs parked)
%   \cmLdCCadLstmDM: an LSTM step of the campaign (LSTMs parked)
%   \cmLdCSeq: an LSTM step of the campaign (LSTMs parked)
%   \cmLdCSeqQ: an LSTM step of the campaign (LSTMs parked)
%   \cmLdCSeqT: an LSTM step of the campaign (LSTMs parked)
%   \cmLdCSeqDM: an LSTM step of the campaign (LSTMs parked)
%   \cmPNewLive: live_feasible input set (live-feasible removed from the paper)
%   \cmPOldLive: live_feasible input set (live-feasible removed from the paper)
%   \cmKeptLive: live_feasible input set (live-feasible removed from the paper)
%   \cmKeptLiveMin: live_feasible input set (live-feasible removed from the paper)
%   \cmKeptLiveMax: live_feasible input set (live-feasible removed from the paper)
%   \cmLstmCadLive: live_feasible input set (live-feasible removed from the paper)
%   \cmLstmCadBase: LSTM (LSTMs parked)
%   \cmLstmCadAll: LSTM (LSTMs parked)
%   \cmLstmCadBelow: LSTM (LSTMs parked)
%   \cmLstmCadBelowSig: LSTM (LSTMs parked)
%   \cmLstmVsRidgeMin: LSTM (LSTMs parked)
%   \cmLstmVsRidgeMax: LSTM (LSTMs parked)
%   \cmIntraVsRidgeMin: LSTM (LSTMs parked)
%   \cmIntraVsRidgeMax: LSTM (LSTMs parked)
%   \cmIntraAllQ: LSTM (LSTMs parked)
%   \cmIntraAllDM: LSTM (LSTMs parked)
%   \cmNlBestVsH: against the former headline (no headline model)
%   \cmNlBestVsHT: against the former headline (no headline model)
%   \cmNlVsHAbove: against the former headline (no headline model)
%   \cmNlVsHAboveSig: against the former headline (no headline model)
%   \cmPnlTotal: D.14, the former headline's P&L decomposition (parked, user decision)
%   \cmPnlTotalX: D.14, the former headline's P&L decomposition (parked, user decision)
%   \cmPnlTopTenShare: D.14, the former headline's P&L decomposition (parked, user decision)
%   \cmPnlTopTwentyShare: D.14, the former headline's P&L decomposition (parked, user decision)
%   \cmPnlTopTwentyShareX: D.14, the former headline's P&L decomposition (parked, user decision)
%   \cmPnlExTwenty: D.14, the former headline's P&L decomposition (parked, user decision)
%   \cmPnlExTwentySh: D.14, the former headline's P&L decomposition (parked, user decision)
%   \cmPnlMktTwenty: D.14, the former headline's P&L decomposition (parked, user decision)
%   \cmPnlMktTwentyExp: D.14, the former headline's P&L decomposition (parked, user decision)
%   \cmPnlMktTwentyP: D.14, the former headline's P&L decomposition (parked, user decision)
%   \cmPnlDaysAll: D.14, the former headline's P&L decomposition (parked, user decision)
%   \cmPnlDaysAllPct: D.14, the former headline's P&L decomposition (parked, user decision)
%   \cmPnlMEn: D.14, the former headline's P&L decomposition (parked, user decision)
%   \cmPnlMEbuys: D.14, the former headline's P&L decomposition (parked, user decision)
%   \cmPnlME: D.14, the former headline's P&L decomposition (parked, user decision)
%   \cmPnlMEShort: D.14, the former headline's P&L decomposition (parked, user decision)
%   \cmPnlNotMEn: D.14, the former headline's P&L decomposition (parked, user decision)
%   \cmPnlNotME: D.14, the former headline's P&L decomposition (parked, user decision)
%   \cmPnlDiffBuyShare: D.14, the former headline's P&L decomposition (parked, user decision)
%   \cmPnlDiff: D.14, the former headline's P&L decomposition (parked, user decision)
%   \cmPnlDiffX: D.14, the former headline's P&L decomposition (parked, user decision)
%   \cmPnlDiffExTen: D.14, the former headline's P&L decomposition (parked, user decision)
%   \cmPnlDiffHighVix: D.14, the former headline's P&L decomposition (parked, user decision)
%   \cmPnlHighVixN: D.14, the former headline's P&L decomposition (parked, user decision)
%   \cmSvAbove: strategy variations under the former headline (no headline model)
%   \cmSvAboveSig: strategy variations under the former headline (no headline model)
%   \cmSvBestMid: strategy variations under the former headline (no headline model)
%   \cmSvBestMidSh: strategy variations under the former headline (no headline model)
%   \cmSvBestMidT: strategy variations under the former headline (no headline model)
%   \cmSvWingN: strategy variations under the former headline (no headline model)
%   \cmSvWingMid: strategy variations under the former headline (no headline model)
%   \cmSvWingX: strategy variations under the former headline (no headline model)
%   \cmSvWingMidSig: strategy variations under the former headline (no headline model)
%   \cmSvWingXSig: strategy variations under the former headline (no headline model)
%   \cmPtFiredRidgeLean: 2026-10-08 parked (lean set; replaced by an MCS)
%   \cmPtPenRidgeLean: 2026-10-08 parked (lean set; replaced by an MCS)
%   \cmPtFiredLassoLean: 2026-10-08 parked (lean set; replaced by an MCS)
%   \cmPtPenLassoLean: 2026-10-08 parked (lean set; replaced by an MCS)
%   \cmPtFiredEnetLean: 2026-10-08 parked (lean set; replaced by an MCS)
%   \cmPtPenEnetLean: 2026-10-08 parked (lean set; replaced by an MCS)
%
% statements of the text checked on these numbers (2026-10-06; not asserted):
%   holds   5.4: sign(s) lifts every one of the paper's eight forecasts above always short (mid)
%   holds   5.4: R trades above always short (cmRefVsShort > 0)
%   holds   5.4: the FOMC columns' Sharpe gap is not significant (HAC t)
%   holds   4.5: every forecast of the model set leans short (sells on > half the days)
%   holds   5.4: the model set's top Sharpe (subpt_enet_all_features) is not significantly above R (HAC t)
%   holds   5.4: no step of Table 8 moves the Sharpe ratio significantly (HAC t; 0 do)
%   holds   5.4: daily refits make LightGBM more accurate (cmLdCCadLgbmQ < 0)
%   holds   5.4: ... without moving its trade significantly (HAC t)
%   holds   5.4: tuning does not help (TuneRs: QLIKE % > 0)
%   holds   5.4: tuning does not help (TuneOp: QLIKE % > 0)
%   holds   5.4: the de-dup moves the per-bar ridge's QLIKE only at rounding
%   holds   5.4: the de-dup leaves the trade
%   holds   5.4: no position of the per-bar ridge changed
%   FLIPPED 5.4: the specification step moves little on the trade (cmLdSpec)
%   holds   5.4: per-bar beats pooled on QLIKE on HAR + calendar (cmLdPerBarBaseQ < 0)
%   holds   5.4: ... significantly (DM)
%   holds   5.4: ... but not on the trade (HAC t not significant)
%   holds   5.4: no linear forecast's trade moved significantly with the de-dup (HAC t)
%   holds   5.4: the mask made no tree significantly better (DM / HAC t)
%   FLIPPED 5.4: no tree more accurate than the ridge (2 QLIKE differences below 0)
%   FLIPPED 5.4: no tree trades above the ridge (7 Sharpe differences above 0)
%   FLIPPED 5.4: daily refits lower every tree's QLIKE
%   holds   5.4: daily refits move XGBoost on all_features' trade the most
%   FLIPPED 5.4 (was): that move is resolved (HAC t significant)
%   holds   5.4: random search is never more accurate than the shipped settings (0 below)
%   holds   5.4: Optuna is never more accurate than the shipped settings (0 below)
%   holds   5.4: no tuned tree's trade differs significantly from the shipped one (0)
%   FLIPPED 5.4: tuning more often is never less accurate (6 above)
%   holds   5.4: ... and never significantly different on the trade (HAC t)
%   FLIPPED 5.4: the best of 50 trials is never more accurate than of 10 / 25 (7 below)
%   FLIPPED 5.4: at crossed fills the straddle ranks first for every forecast of the model set
%   holds   5.4: wings cost at every width at crossed fills, for every forecast
%   holds   5.4: for every forecast the 20 best days are buys (the tail-trade reading)
%   FLIPPED 4.2: the B contamination moves no Sharpe ratio by 0.1 (model set)
%   FLIPPED 4.5: dropping the censored days moves no Sharpe ratio by more than 0.06
%   holds   4.x: the margin rule lets few retunes fire
%   FLIPPED methods: pre-tuning improves the pooled models' QLIKE (mean change < 0)
\fi

% 2026-10-15 (user): parked -- section renamed ("experiments + results"); label unchanged. Was: \section{Results}
\section{Experiments and results}
\label{sec:rh_results}
% 2026-10-22 (user): rerun note.
% 2026-10-23 (user, final rerun): parked -- the note before the window choice; the window is decided and the generators read the final rerun (writeup/final_run.py). Was: "\editc{orange}{RR}{Decision: a rerun of the pooled intraday models is approved (22 October 2026) and waits for the window choice.\\Problem: the pooled rows of Table~\ref{tab:rh_oos} and every pooled number in this section come from forecasts with a 1,000-day window that predate the late-start, no-observation and KKT-check fixes and still include the FOMC statement dummy; they will move. See the notes to the authors, ``Rerun pending''.}"
% 2026-10-23 (user, final rerun, phase 2): parked -- note RR resolved: the generators read the final rerun (REHAUL_FINAL_RUN=1), the window robustness is in Table tab:rh_robust_w3500 (item "Robustness to the window length"), and the subperiod directions and the signs of the gaps to the closing-interval HAR benchmark were rechecked on the final values; the MinGRU stays parked. Was: "\editc{orange}{RR}{Decision: the pooled intraday models are rerun on the final design (decided 23 October 2026; writeup/final\_run.py): first evaluation day \rhorSpanStart, the linear penalties chosen once on \rhorPenSample{} and frozen, a rolling window of \rhorWindowDays{} trading days, re-estimated at every bar, with the late-start, no-observation and KKT-check fixes and without the FOMC statement dummy. Table~\ref{tab:rh_oos} and the pooled numbers of this section are generated from those forecasts (REHAUL\_FINAL\_RUN=1).\\Pending (2026-10-23): the robustness table at the longer window (the random forest's walk at that window is not finished); the sentences whose wording depends on the final values (the subperiod comparison and the signs of the gaps to the closing-interval HAR benchmark below) are rechecked in the second pass. The MinGRU stays parked.}"
% 2026-10-17 (user): prose rewritten in the style of Zhang, Zhang, Cucuringu, and Qian (2024).
\subsection{Experimental setup}
\label{sec:rh_exp_setup}
% 2026-10-15 (user): the full implementation material ((a) samples and windows, (b) training schemes in practice,
% (c) tuned hyperparameters, (d) tuning procedures; labels sec:rh_exp_samples, sec:rh_exp_schemes,
% sec:rh_methods_hyper, sec:rh_methods_tuning kept) moved to the appendix, "Implementation details" (app:rh_impl,
% 10_appendix.tex), with its tables (tab:rh_dm_tuning, tab:rh_dm_hyper, tab:rh_dm_hyper_search; 10_appendix_floats.tex).
% This subsection keeps a short summary.
\begin{itemize}
% 2026-10-15 (user): parked -- "Out of sample period doesn't make sense for forecasts because we're performing rolling
% backtests. The forecasting results need to be calculated to reflect such (that the out of sample forecasts are done
% on a rolling basis after the burn in training window)". Was:
% \item \textbf{Out-of-sample period} $\mathcal T_{\mathrm{oos}}$: \rheSampleStart{} to \rheSampleEnd,
%   $T_{\mathrm{oos}} = |\mathcal T_{\mathrm{oos}}| = \cmNdays$ trading days
%   (Section~\ref{sec:rh_data}).
% 2026-10-22 (user): moved here from the data section (03_data.tex, "Estimation and evaluation samples"): the pooled
% intraday sample and the closing-interval sample; the evaluation span below was already here.
\item \textbf{Estimation samples} (see Appendix Table~\ref{tab:rh_dm_samples}). The pooled intraday sample comprises
  \rhdmPooledRows{} half-hour intervals beginning in June 1998,
  \[
    \rhdmPooledRows = \rhdmGridRows \text{ intervals on the grid} - \rhdmBurnIn \text{ burn-in samples};
  \]
  the pooled linear models and tree ensembles of Table~\ref{tab:rh_oos} and of this section and
  Section~\ref{sec:rh_econ} are trained on its \rhorNBars{} regular-hours bars with rolling windows of
  $W = \rhorWindowBars$ bars (\rhorWindowDays{} trading days; training scheme below), and the pooled linear runs of the
  model set on all its intervals with $W = \rhdmPooledLinWindow$ intervals (Appendix~\ref{sec:rh_exp_schemes}).
  % 2026-10-22 (user): the closing-interval sample is not presented in the experiments section; it is described with
  % the closing-interval models in the final subsection of Section 7 (sec:rh_econ_last30) and in
  % Appendix app:rh_closing_sample.
% 2026-10-22 (user): parked -- the closing-interval pre-sample (\cmPt... macros) is not presented in the experiments
% section (see the final subsection of Section 7). Was: "This evaluation span begins with the pooled intraday linear
% forecasts immediately following the pre-sample validation period (\cmPtValFirst{} to \cmPtValLast) used to select the
% linear regularization penalties."
\item \textbf{Evaluation span $\mathcal T_{\mathrm{oos}}$} (see Section~\ref{sec:rh_results_eval}). We evaluate
  the rolling out-of-sample forecasts over the period from \rhoosSpanStart{} to \rhoosSpanEnd, which comprises
  $T_{\mathrm{oos}} = |\mathcal T_{\mathrm{oos}}| = \rhoosDays$ trading days on which all competing forecasts are
  available (after excluding \rhoosEarlyCloses{} early-close sessions).
  % 2026-10-20 (user): added -- Table 1 scored from 2012 (TABLE1_SWITCH.md, approved).
  % 2026-10-22 (user, E36): parked -- Table 1 now holds the pooled tree ensembles too. Was: "The pooled linear models of Table~\ref{tab:rh_oos} are scored on a longer span, from \rhorSpanStart{} to \rhorSpanEnd{} (\rhorDays{} trading days, after excluding \rhorEarlyCloses{} early-close sessions), with a rolling window of \rhorWindowDays{} trading days and penalties chosen once on \rhorPenFolds{} validation folds covering \rhorPenYears{}, before the first forecast, and then frozen."
  % 2026-10-23 (final-run switch): parked -- "\rhorPenFolds{} validation folds covering \rhorPenYears{}" prints "1 validation folds" on the final rerun (one pre-sample fold); \rhorPenSample is the phrase for either run (make_rehaul_oos_rerun.py). Was: "the linear penalties are chosen once on \rhorPenFolds{} validation folds covering \rhorPenYears{}, before the first forecast, and then frozen,"
  The pooled linear models and tree ensembles of Table~\ref{tab:rh_oos} are scored on a longer span, from
  \rhorSpanStart{} to \rhorSpanEnd{} (\rhorDays{} trading days, after excluding \rhorEarlyCloses{} early-close
  sessions), with a rolling window of \rhorWindowDays{} trading days; the linear penalties are chosen once on
  \rhorPenSample, before the first forecast, and then frozen, and each
  tree ensemble keeps one configuration from an Optuna search on the pre-sample.
% 2026-10-16 (user): parked -- option-data subsample removed from Experimental setup. The set is defined in Data
% (sec:rh_data); the other \rhoosNonOptDays{} days enter forecast accuracy only, now in Economic value
% (sec:rh_econ_measures). Was:
% \item \textbf{Option-data subsample} $\mathcal T_{\mathrm{opt}} \subset \mathcal T_{\mathrm{oos}}$: the
%   $T_{\mathrm{opt}} = \rhrDays$ days with a same-day SPXW expiration (Section~\ref{sec:rh_data}), on which the straddle
%   economics of Section~\ref{sec:rh_econ} are evaluated; the other \rhoosNonOptDays{} days enter forecast accuracy only.
% 2026-10-15 (user): parked -- last-30-min scheme moved to the end of the economics section (sec:rh_econ_last30);
% the pooled intraday scheme is the training scheme of this section, and the last-30-minutes scheme keeps one line
% (the benchmark forecast, below). Was:
% \item \textbf{Closing-interval scheme:} the last 30 minutes of each trading day; $W = \cmPtFitWin$ trading days.
% \item \textbf{Pooled intraday scheme:} all $M = \rhdmNIntervals$ intervals of each day; $W = \rhdmPooledWindow$
%   intervals for the trees and $\rhdmPooledLinWindow$ for the linear models.
% 2026-10-22 (user, E34): parked -- the tree ensembles are the base-5 refit on the regular-hours bars (writeup/pooled_trees_13.py). Was: "\item \textbf{Training scheme.} Under the pooled intraday estimation scheme, models are fit across all $M = \rhdmNIntervals$ half-hour intervals of each trading day using a rolling window of $W = \rhdmPooledWindow$ intervals for the tree ensembles and $W = \rhdmPooledLinWindow$ intervals for the linear specifications. The linear models of Table~\ref{tab:rh_oos} are fit on the \rhorNBars{} regular-hours bars with a rolling window of \rhorWindowBars{} bars (\rhorWindowDays{} trading days). On the full predictor set, the pooled LightGBM and XGBoost forecasts reflect an earlier specification and the pooled random forest forecast is preliminary (see Table~\ref{tab:rh_oos})."
% 2026-10-22 (user, E34: one design for the 12 pooled forecasts): parked -- the two designs; the pooled linear models of this section and of Section 7 are the forecasts of Table 1 on the trees' design. Was: "\item \textbf{Training scheme.} Under the pooled intraday estimation scheme, the linear specifications are fit across all $M = \rhdmNIntervals$ half-hour intervals of each trading day using a rolling window of $W = \rhdmPooledLinWindow$ intervals, and the tree ensembles across the \rheTreeNBars{} regular-hours bars of each trading day using a rolling window of $W = \rheTreeWindowBars$ bars (\rheTreeWindowDays{} trading days), re-estimated at every bar."
\item \textbf{Training scheme.} Under the pooled intraday estimation scheme, the linear specifications and the
  tree ensembles are fit across the \rhePoolNBars{} regular-hours bars of each trading day using a rolling window
  of $W = \rhePoolWindowBars$ bars (\rhePoolWindowDays{} trading days), re-estimated at every bar; the pooled
  linear runs of the model set use all $M = \rhdmNIntervals$ half-hour intervals with $W = \rhdmPooledLinWindow$
  intervals.
  % 2026-10-22 (user, E36): parked -- Table 1 now holds the tree ensembles too. Was: "The linear models of Table~\ref{tab:rh_oos} are fit on the \rhorNBars{} regular-hours bars with a rolling window of \rhorWindowBars{} bars (\rhorWindowDays{} trading days), the design of the tree ensembles."
  The models of Table~\ref{tab:rh_oos}, the linear models and the tree ensembles, are fit on the \rhorNBars{}
  regular-hours bars with a rolling window of \rhorWindowBars{} bars (\rhorWindowDays{} trading days).
  % 2026-10-16 (user): parked -- attach each accuracy row to the configuration that produced it.
  % The hyperparameter rule is now one frozen pre-sample pick. The sample and the window stay open;
  % the comparison is in the author notes. Was:
  % \editc{orange}{E15}{Decision: attach each accuracy row to the configuration that produced it. Section~\ref{sec:rh_exp_schemes} says the intraday trees use one configuration, never re-selected, re-estimated every interval, on $W=\rhdmPooledWindow$. Table~\ref{tab:rh_dm_hyper} prints that configuration (LightGBM $J$, $\eta$, $B$, $n_{\min}$; XGBoost depth, $\eta$, $B$; the random-forest trial). This bullet says the full-predictor-set LightGBM and XGBoost rows of Table~\ref{tab:rh_oos} are an earlier draft and the random forest is preliminary. Note (a) of that table also marks the intraday OLS on the HAR benchmark. The earlier draft's $J$, $\eta$, $B$, and $n_{\min}$ are not in Table~\ref{tab:rh_dm_hyper}, so those QLIKE cells are not the losses of the printed configuration.}
  % 2026-10-22 (user, E34): parked -- the tree clauses: the pooled trees are now one frozen configuration per method refit at every bar on the 13 regular-hours bars with the 1,000-day window (writeup/pooled_trees_13.py). Was: "...the HAR-benchmark $\lambda$ path, the full-set LightGBM draft, the full-set XGBoost draft, and the preliminary random forest were not produced by one frozen configuration. ... and the full-set LightGBM row, the full-set XGBoost row, and the random-forest row are an earlier draft or a preliminary fit. Those QLIKE numbers are not the losses of one shared configuration, and the table has not been refit."
  % 2026-10-22 (user, E34: one design for the 12 pooled forecasts): parked -- E15's open sample and window and the cells not produced by one frozen configuration; Table 1 is the 2012 rerun with frozen penalties and the trees (E36), and Sections 6-7 use those forecasts. Was: "\editc{orange}{E15}{Decision: one hyperparameter rule for every class, frozen after a single pre-sample pick.\\Problem: the sample and the window are still open, and the printed QLIKE cells for the linear light re-tune and the HAR-benchmark $\lambda$ path were not produced by one frozen configuration.\\ The rule for the comparison is one configuration per model, chosen once on the pre-sample and then held fixed for the whole evaluation span. Two inputs to that rule are still unset, and the only write-up of the comparison is in the author notes. The sample is either every one of the $\rhdmNIntervals$ intervals or the 13 regular-hours bars whose end stamps run from 10:00 to 16:00. The window is either $W=\rhdmPooledWindow$ intervals or $W=\rhdmPooledLinWindow$ intervals. On top of that open choice, some cells already printed in Table~\ref{tab:rh_oos} were not produced by one frozen configuration. The linear penalties still move under the light re-tune and the HAR-benchmark penalties move on the dates left open in note E23. Those QLIKE numbers are not the losses of one shared configuration, and the table has not been refit. The pooled trees are now one frozen configuration per method, refit at every bar on the \rheTreeNBars{} regular-hours bars with a window of \rheTreeWindowDays{} trading days (note E34).}"
  % 2026-10-23 (user, E15, final rerun): parked -- note E15 closed: one frozen pre-sample pick per class is the final design (the linear penalties on pre-sample fold 3, validation year 2013, results/window_sweep_2014start_2026-10-22/w3000/frozen_penalties.json; the trees' frozen Optuna configurations, not re-tuned: results/pooled_trees_final_2014start_w3000/run.json, results/pooled_rf_skx_2014start_w3000/RULE_rf_skx.txt). The rule is restated with macros in the item "Rolling re-estimation with frozen hyperparameters" below. Was: "\editc{orange}{E15}{Decision: one hyperparameter rule for every class, frozen after a single pre-sample pick.\\Problem (2026-10-22): for the pooled models the rule now holds; the note can be closed once confirmed.\\ Every pooled model of Table~\ref{tab:rh_oos} and of Sections~\ref{sec:rh_results}--\ref{sec:rh_econ} is fit on the \rhePoolNBars{} regular-hours bars with a window of \rhePoolWindowDays{} trading days ($W = \rhePoolWindowBars$ bars), re-estimated at every bar: the linear penalties are chosen once on \rhorPenFolds{} pre-sample validation folds covering \rhorPenYears{} and then frozen, and each tree method keeps one Optuna configuration (notes E34, E36). The tree configurations were selected on $W=\rhdmPooledWindow$ intervals and are used as selected (note E24). The pooled linear runs of the model set on all $\rhdmNIntervals$ intervals ($W=\rhdmPooledLinWindow$) are no longer the pooled rows of these sections.}"
  % 2026-10-15 (finish): FLAG resolved -- typed 96{,}000 -> \rhdmPooledLinWindow (make_rehaul_data_methods.py, same value).
  % 2026-10-15 (user, restructure): the flags of Table tab:rh_oos (note a, dagger) restated, as these forecasts now
  % carry the main results (make_rehaul_results.py: EARLIER_POOLED_FULL, INTERIM).
% 2026-10-17 (user): parked -- closing-interval material (relabeled: closing-interval HAR benchmark); the window moved to Appendix app:rh_closing_bench. Was: "Our baseline benchmark is the daily HAR regression estimated on a rolling window of $W = \cmPtFitWin$ trading days;"
% 2026-10-17 (user, new-notes 11): parked -- a HAR-type model, not Corsi's daily HAR (src/features/extractors/har.py:18-21). Was: "the closing-interval HAR benchmark, the daily HAR regression estimated on closing-interval observations"
\item \textbf{Benchmark forecast.} Our baseline benchmark is the closing-interval HAR benchmark, a HAR-type model in the
  % 2026-10-18 (user, new-notes 6): parked -- k in samples. Was: "spirit of \citet{Corsi2009} on the target's trailing averages over $k \in \{1, 5, 25, 125, 625, 3125\}$ 30-minute"
  spirit of \citet{Corsi2009} on the target's trailing averages over $k \in \{1, 5, 25, 125, 625, 3125\}$ samples
  % 2026-10-18 (user, new-notes 6): parked -- k in samples. Was: "grid intervals and the \rhdmNCalendar{} calendar columns, estimated by OLS on closing-interval observations"
  and the \rhdmNCalendar{} calendar columns, estimated by OLS on closing-interval observations
  (Appendix~\ref{app:rh_closing_bench});
  % 2026-10-17 (user): parked -- closing-interval material moved to the appendix; Section -> Appendix. Was: "window of $W = \cmPtFitWin$ trading days; Section~\ref{sec:rh_econ_last30} examines all model classes trained"
  Appendix~\ref{sec:rh_econ_last30} examines all model classes trained
  analogously on the closing 30-minute interval alone.
% 2026-10-17 (user): parked -- closing-interval material (the daily refit is in Appendix app:rh_closing_bench). Was: "Models are re-estimated at each period (see"
% 2026-10-20 (user, third-list 1): parked -- tuning wording matched to the pick-once-and-freeze decision (E23; Appendices B.2 and D). Was: "\item \textbf{Rolling re-estimation and causal tuning.} The pooled intraday models are re-estimated at every 30-minute interval (see Appendix~\ref{sec:rh_exp_schemes}), with hyperparameters either held fixed or selected causally on an initial pre-sample or trailing validation windows prior to the forecast origin."
% 2026-10-23 (user, E15 closed, final rerun): the rule of note E15 restated with macros (which pre-sample, which
% configuration). Was: "Their hyperparameters are chosen once, on the fixed pre-sample, and then held fixed for the whole
% evaluation span, with no re-tuning (Appendix~\ref{app:rh_hyper})."
\item \textbf{Rolling re-estimation with frozen hyperparameters.} The pooled intraday models are re-estimated at every
  30-minute interval (see
  Appendix~\ref{sec:rh_exp_schemes}). Their hyperparameters are chosen once, before the first forecast, and then
  held fixed for the whole evaluation span, with no re-tuning (Appendix~\ref{app:rh_hyper}): the linear penalties
  on \rhorPenSample, and for each tree ensemble the one configuration that an Optuna search selected on the
  pre-sample with a window of \rhdmPooledWindow{} intervals (Appendix~\ref{sec:rh_methods_hyper}).
  % 2026-10-23 (user, final rerun, phase 2): added -- the HAR-set penalties of the final rerun
  % (results/window_sweep_2014start_2026-10-22/w3000/frozen_penalties.json = \rhdmPlPen...Base); checked: the HAR-set
  % ridge, lasso and elastic-net forecasts differ from the OLS forecast by at most 5e-6 in relative terms (forecasts.parquet),
  % and their rows in Table 1 equal the OLS row.
  On the HAR benchmark the frozen penalties at this window are \rhdmPlPenRidgeBase{} (ridge), \rhdmPlPenLassoBase{}
  (lasso) and \rhdmPlPenEnetBase{} (elastic net), numerically zero: the three penalized models reproduce the OLS
  forecast, and their rows in Table~\ref{tab:rh_oos} coincide with the OLS row.
% 2026-10-15 (user): parked -- the lags without their samples. Was: "Diebold--Mariano tests with the
% Harvey--Leybourne--Newbold correction ($q_{\mathrm{NW}} = \fqsDMLag$) and HAC $t$-statistics ($q_{\mathrm{NW}} =
% \fqsHACLag$), each with ..."
% 2026-10-16 (user): parked -- the straddle HAC clause. The lag on $\mathcal T_{\mathrm{opt}}$ is in Economic value
% (sec:rh_econ_measures). Was: "... and HAC $t$-statistics of straddle returns ($q_{\mathrm{NW}} = \fqsHACLag$ on
% $\mathcal T_{\mathrm{opt}}$), each with a Newey--West long-run variance ..."
\item \textbf{Evaluation metric and inference.} We evaluate predictive accuracy via the QLIKE loss against actual
  realized variance (see Section~\ref{sec:rh_results_eval}) and conduct pairwise inference using Diebold--Mariano
  tests with the Harvey--Leybourne--Newbold small-sample adjustment and a Newey--West long-run variance estimator
  ($q_{\mathrm{NW}} = \rhoosDMLag$ on
  $\mathcal T_{\mathrm{oos}}$).\footnote{Appendix~\ref{sec:rh_methods_tests} provides the complete specification
  of these statistical tests.}
  % 2026-10-16 (user): the parenthetical section pointer became the footnote above. Was:
  % "(Section~\ref{sec:rh_methods_tests})."
  % 2026-10-15 (user): the lag values moved here from Methods (Forecast evaluation and Statistical tests, numbers
  % removed there); pointer to Forecast evaluation added. Was: "QLIKE against the actual realized variance;
  % Diebold--Mariano tests with the Harvey--Leybourne--Newbold correction and HAC $t$-statistics
  % (Section~\ref{sec:rh_methods_tests})."
\end{itemize}
We refer the reader to Appendix~\ref{app:rh_impl} for comprehensive implementation details.

% 2026-10-15 (user): new subsection -- the out-of-sample accuracy table (tab:rh_oos), generator
% writeup/make_rehaul_oos_table.py (macros \rhoos... in generated/rehaul_oos_numbers.tex).
\subsection{Out-of-sample accuracy}
\label{sec:rh_results_oos}
\begin{itemize}
% 2026-10-15 (user): parked -- the option-data days; accuracy over every rolling forecast. Was: "... estimation
% samples, over the same \rhrDays{} trading days."
% 2026-10-15 (user): parked -- last-30-min scheme moved to the end of the economics section; the table (unchanged for
% now) is described by its pooled intraday columns, and its "Last 30 minutes" columns are discussed in
% sec:rh_econ_last30. Was: "\textbf{Overview:} Table~\ref{tab:rh_oos} reports the out-of-sample MSE, average QLIKE and
% $R^2_{\mathrm{OOS}}$ of every model on the HAR benchmark and on all features, for the closing-interval (last 30
% minutes) and the pooled intraday estimation samples, over the same \rhoosDays{} trading days of the evaluation span
% $\mathcal T_{\mathrm{oos}}$."
\item \textbf{Overview of forecasting performance.} Table~\ref{tab:rh_oos} reports the out-of-sample MSE, average
  QLIKE loss, and out-of-sample $R^2_{\mathrm{OOS}}$ for each pooled intraday model class under both the HAR
  % 2026-10-20 (user): parked -- Table 1 scored from 2012 (TABLE1_SWITCH.md, approved). Was: "benchmark and the full predictor set across the \rhoosDays{} trading days of the evaluation span $\mathcal T_{\mathrm{oos}}$. The corresponding metrics for specifications estimated exclusively on the closing"
  benchmark and the full predictor set across the \rhorDays{} trading days from \rhorSpanStart{} to
  \rhorSpanEnd. The corresponding metrics for specifications estimated exclusively on the closing
  % 2026-10-17 (user): parked -- closing-interval material (relabeled: closing-interval HAR benchmark). Was: "including the daily OLS regression on the HAR benchmark that serves as our benchmark"
  % 2026-10-18 (user, new-notes 10): parked -- closing cadence only in the closing section. Was: "half-hour interval—including the daily OLS regression on the HAR benchmark predictor set, our closing-interval HAR benchmark"
  half-hour interval—including the OLS regression on the HAR benchmark predictor set, our closing-interval HAR benchmark
  % 2026-10-17 (user): parked -- closing-interval material moved to the appendix; Section -> Appendix. Was: "forecast—are presented in Table~\ref{tab:rh_oos_last30} (see Section~\ref{sec:rh_econ_last30})."
  forecast—are presented in Table~\ref{tab:rh_oos_last30} (see Appendix~\ref{sec:rh_econ_last30}).
  % 2026-10-20 (user): added -- "add the one-sentence 2012-2019 / 2020-2024 split in the text (no table)"; numbers from make_rehaul_oos_rerun.py (\rhorSub...).
  % 2026-10-20 (user): parked -- the years of the split as macros (no typed numbers). Was: the same sentence with the years typed ("Split at the start of 2020", "2012--2019", "2020--2024").
  % 2026-10-22 (user, E36): parked -- the split macros cover the linear rows only, and Table 1 now holds the trees. Was: "Split at the start of \rhorSubBStartYear, the full predictor set has the lower QLIKE in both parts"
  For the linear models, split at the start of \rhorSubBStartYear, the full predictor set has the lower QLIKE in both parts
  (\rhorSubALabel: \rhorSubAQlikeAfMin--\rhorSubAQlikeAfMax{} against \rhorSubAQlikeHar{} for the HAR benchmark;
  \rhorSubBLabel: \rhorSubBQlikeAfMin--\rhorSubBQlikeAfMax{} against \rhorSubBQlikeHar), and a higher
  $R^2_{\mathrm{OOS}}$ on \rhorSubALabel{} (\rhorSubARoosAfMin--\rhorSubARoosAfMax\% against \rhorSubARoosHar\%)
  but a lower one on \rhorSubBLabel{} (\rhorSubBRoosAfMin--\rhorSubBRoosAfMax\% against \rhorSubBRoosHar\%), so
  its lower $R^2_{\mathrm{OOS}}$ over the whole span comes from \rhorSubBLabel.
  % 2026-10-15 (user): Table 1 split ("Split it"); was "(columns ``Pooled intraday'') ... The columns ``Last 30
  % minutes'' are discussed in Section~\ref{sec:rh_econ_last30}; their OLS row is the benchmark forecast."
  % 2026-10-20 (user): "we need to straighten this out so it's like for like"; "evaluate on all 48 periods".
  % 2026-10-20 (user): parked -- "just keep the 13 session panel on it": E36 named both panels. Was the same note with Panel A (48 intervals) and Panel B (13 bars) and the tests of both.
  % 2026-10-20 (user): parked -- "switch the main window to 1,000 days, not 2,000. Table 1 should cover 2012-2024 as a single panel". Was: "Decision (user, 2026-10-20): Table~\ref{tab:rh_oos} is like for like. It holds the pooled linear rerun trained and scored on the 13 regular-hours bars (base-5 HAR horizons, FOMC series in, penalties chosen once on the pre-sample, a 2{,}000-day window), with every bar scored. Its tests are Diebold--Mariano on the daily average QLIKE. The 48-interval rerun is not in the table; its numbers are in results/oos\_rerun\_2026-10-20/."
  % 2026-10-20 (user): parked -- the prose now names both spans (TABLE1_SWITCH.md applied). Was: "\editc{orange}{E36}{Decision (user, 2026-10-20): Table~\ref{tab:rh_oos} is like for like. It holds the pooled linear rerun trained and scored on the 13 regular-hours bars (base-5 HAR horizons, FOMC series in, a rolling window of \rhorWindowDays{} trading days, penalties chosen once on four validation folds covering \rhorPenYears{} and then frozen), with every bar scored from \rhorSpanStart{} to \rhorSpanEnd{} (\rhorDays{} trading days, one panel). Its tests are Diebold--Mariano on the daily average QLIKE. The 2{,}000-day reruns (from 2020 and from 2012) and the 48-interval rerun are not in the table; their numbers are in results/oos\_rerun\_2026-10-20/ and results/oos\_span\_check\_2026-10-20/. The rest of this section, the abstract, the introduction and the conclusion still describe the span from \rhoosSpanStart{}; results/oos\_span\_check\_2026-10-20/TABLE1\_SWITCH.md proposes the changes. Tree ensembles are held off until they are refit on the cluster.\\Problem: this bullet and the $R^2_{\mathrm{OOS}}$ bullet below still describe the earlier table, and so do the abstract, the introduction and the conclusion. They speak of the 15:30 forecast of the closing half hour, six model classes with trees, a closing-half-hour $R^2$ benchmark, and the \texttt{\textbackslash rhrAfPool...} counts that writeup/make\_rehaul\_results.py computed from the earlier runs.\\New tests, QLIKE, all features against the HAR benchmark of the same model, $p$-values in parentheses: ridge \rhorDmRidgeB{} (\rhorPRidgeB), lasso \rhorDmLassoB{} (\rhorPLassoB), elastic net \rhorDmEnetB{} (\rhorPEnetB). The like-for-like values of the prose numbers are being computed (branch insurance-selling-2026-09-19-prose-numbers).}"
  % 2026-10-20 (user): parked -- the R2_OOS bullet and the appendix scoring line are fixed. Was: "\editc{orange}{E36}{Decision (user, 2026-10-20): Table~\ref{tab:rh_oos} is like for like. It holds the pooled linear rerun trained and scored on the \rhorNBars{} regular-hours bars (the HAR horizons of Section~\ref{sec:rh_predictors}, FOMC series in, a rolling window of \rhorWindowDays{} trading days, penalties chosen once on \rhorPenFolds{} validation folds covering \rhorPenYears{} and then frozen), with every bar scored from \rhorSpanStart{} to \rhorSpanEnd{} (\rhorDays{} trading days, one panel). Its tests are Diebold--Mariano on the daily average QLIKE. The reruns with the longer window and the rerun on every interval of the day are not in the table; their numbers are in results/oos\_rerun\_2026-10-20/ and results/oos\_span\_check\_2026-10-20/. The abstract, the introduction, Data, this section, the conclusion and the appendix now give both spans: Table~\ref{tab:rh_oos} from \rhorSpanStart, the closing-interval models, Table~\ref{tab:rh_oos_last30} and the straddle sample from \rhoosSpanStart. Tree ensembles are held off until they are refit on the cluster.\\Problem: the $R^2_{\mathrm{OOS}}$ bullet below, the abstract, the introduction and the conclusion still carry the earlier table's framing: the 15:30 forecast of the closing half hour, six model classes with trees, a closing-half-hour $R^2$ benchmark, and the \texttt{\textbackslash rhrAfPool...} counts that writeup/make\_rehaul\_results.py computed from the earlier runs. The Table~\ref{tab:rh_dm_samples} note (writeup/make\_rehaul\_data\_methods.py) still gives the linear window as $W = \rhdmPooledLinWindow$ intervals.\\New tests, QLIKE, all features against the HAR benchmark of the same model, $p$-values in parentheses: ridge \rhorDmRidgeB{} (\rhorPRidgeB), lasso \rhorDmLassoB{} (\rhorPLassoB), elastic net \rhorDmEnetB{} (\rhorPEnetB). The like-for-like values of the prose numbers are being computed (branch insurance-selling-2026-09-19-prose-numbers).}"
  % 2026-10-22 (user, E34): added -- the trees are refit (the sentence after "held off until they are refit on the cluster").
  % 2026-10-22 (user, E36): parked -- the trees enter Table 1, scored on every bar like the linear rows (writeup/make_rehaul_oos_rerun.py). Was: "It holds the pooled linear rerun trained and scored on the \rhorNBars{} regular-hours bars"; "Tree ensembles are held off until they are refit on the cluster. They are now refit on this design (2026-10-22, note E34: results/pooled\_trees\_rerun13\_2026-10-20/w1000, every bar from \rheTreeRunStart); the prose numbers read their 16:00 forecasts, and the table does not hold them yet."; "New tests, QLIKE, all features against the HAR benchmark of the same model, $p$-values in parentheses: ridge \rhorDmRidgeB{} (\rhorPRidgeB), lasso \rhorDmLassoB{} (\rhorPLassoB), elastic net \rhorDmEnetB{} (\rhorPEnetB)."
  % 2026-10-22 (user, E34: one design for the 12 pooled forecasts): parked -- in note E36, "the \texttt{\textbackslash rhrAfPool...} counts that writeup/make\_rehaul\_results.py computed from the earlier runs. The Table~\ref{tab:rh_dm_samples} note (writeup/make\_rehaul\_data\_methods.py) still gives the linear window as $W = \rhdmPooledLinWindow$ intervals."
  % 2026-10-23 (user, E36, final rerun): parked -- the note on the 2012 design (the 1,000-day runs from 3 January 2012, results/pooled_trees_rerun13_2026-10-20/w1000; "\rhorPenFolds{} validation folds covering \rhorPenYears{}" printed "1 validation folds" on the final rerun) and its Problem, now resolved (the abstract, the introduction and the conclusion give both spans and print counts). Was: "\editc{orange}{E36}{Decision (user, 2026-10-20): Table~\ref{tab:rh_oos} is like for like. It holds the pooled linear rerun, and since 2026-10-22 the pooled tree ensembles fitted on its design, trained and scored on the \rhorNBars{} regular-hours bars (the HAR horizons of Section~\ref{sec:rh_predictors}, FOMC series in, a rolling window of \rhorWindowDays{} trading days, penalties chosen once on \rhorPenFolds{} validation folds covering \rhorPenYears{} and then frozen), with every bar scored from \rhorSpanStart{} to \rhorSpanEnd{} (\rhorDays{} trading days, one panel). Its tests are Diebold--Mariano on the daily average QLIKE. The reruns with the longer window and the rerun on every interval of the day are not in the table; their numbers are in results/oos\_rerun\_2026-10-20/ and results/oos\_span\_check\_2026-10-20/. The abstract, the introduction, Data, this section, the conclusion and the appendix now give both spans: Table~\ref{tab:rh_oos} from \rhorSpanStart, the closing-interval models, Table~\ref{tab:rh_oos_last30} and the straddle sample from \rhoosSpanStart. Tree ensembles were held off until they were refit on the cluster. They are now refit on this design (2026-10-22, note E34: results/pooled\_trees\_rerun13\_2026-10-20/w1000, every bar from \rheTreeRunStart), and the table holds them: LightGBM, XGBoost and the random forest on both predictor sets, one configuration per method from an Optuna search on the pre-sample, frozen, scored on the same \rhorNIntB{} bars as the linear rows (asserted) and given the same tests. The prose numbers of Section~\ref{sec:rh_results_predictors} read their 16:00 forecasts over $\mathcal T_{\mathrm{oos}}$.\\Problem: the abstract, the introduction and the conclusion still carry the earlier table's framing: the 15:30 forecast of the closing half hour, six model classes with trees, and the \texttt{\textbackslash rhrAfPool...} counts, which writeup/make\_rehaul\_results.py now computes from the 16:00 rows of the forecasts of this table (2026-10-22, note E34). The Table~\ref{tab:rh_dm_samples} note gives these models' window, $W = \rhePoolWindowBars$ bars.\\New tests, QLIKE, all features against the HAR benchmark of the same model, $p$-values in parentheses: ridge \rhorDmRidgeB{} (\rhorPRidgeB), lasso \rhorDmLassoB{} (\rhorPLassoB), elastic net \rhorDmEnetB{} (\rhorPEnetB), LightGBM \rhorDmLgbmB{} (\rhorPLgbmB), XGBoost \rhorDmXgbB{} (\rhorPXgbB), random forest \rhorDmRfB{} (\rhorPRfB). The like-for-like values of the prose numbers are being computed (branch insurance-selling-2026-09-19-prose-numbers).}"
  % 2026-10-23 (user, final rerun, phase 2): parked -- note E36 resolved: Table 1 holds the final rerun (linear models and tree ensembles, one design, every bar, the same bars in every row), its notes give the design, the robustness table at the longer window is in the appendix (tab:rh_robust_w3500), and the subperiod sentence was rechecked (its directions hold on the final values). Was: "\editc{orange}{E36}{Decision (user, 2026-10-20; final rerun 2026-10-23): Table~\ref{tab:rh_oos} is like for like. It holds the pooled linear models and the pooled tree ensembles of the final rerun (writeup/final\_run.py, writeup/make\_rehaul\_oos\_rerun.py with span \rhorSpanStartYear), trained and scored on the \rhorNBars{} regular-hours bars (the HAR horizons of Section~\ref{sec:rh_predictors}, the \rhorNColsAf{} columns of the full predictor set, a rolling window of \rhorWindowDays{} trading days, re-estimated at every bar; the linear penalties chosen once on \rhorPenSample{} and then frozen, one frozen configuration per tree method), with every bar scored from \rhorSpanStart{} to \rhorSpanEnd{} (\rhorDays{} trading days, one panel) and the same \rhorNIntB{} bars in every row (asserted). Its tests are Diebold--Mariano on the daily average QLIKE. The closing-interval models, Table~\ref{tab:rh_oos_last30} and the straddle sample keep the span from \rhoosSpanStart. The prose numbers of Section~\ref{sec:rh_results_predictors} read the 16:00 forecasts of the same runs over $\mathcal T_{\mathrm{oos}}$.\\Pending (2026-10-23): the robustness table at the longer window (the random forest's walk at that window is not finished); the direction of the subperiod sentence above is rechecked in the second pass.\\New tests, QLIKE, all features against the HAR benchmark of the same model, $p$-values in parentheses: ridge \rhorDmRidgeB{} (\rhorPRidgeB), lasso \rhorDmLassoB{} (\rhorPLassoB), elastic net \rhorDmEnetB{} (\rhorPEnetB), LightGBM \rhorDmLgbmB{} (\rhorPLgbmB), XGBoost \rhorDmXgbB{} (\rhorPXgbB), random forest \rhorDmRfB{} (\rhorPRfB).}"
% 2026-10-20 (user): parked -- the benchmark of Table 1 is the mean of the same bar, not of the closing half hour. Was: "\item \textbf{Computation of $R^2_{\mathrm{OOS}}$.} Following \citet{CampbellThompson2008}, the reference forecast in $R^2_{\mathrm{OOS}}$ is the expanding historical mean of closing-half-hour realized variance computed from \rhoosBenchStart{} through day $t - 1$, which relies solely on information available prior to the forecast origin."
\item \textbf{Computation of $R^2_{\mathrm{OOS}}$.} Following \citet{CampbellThompson2008}, the reference
  forecast in $R^2_{\mathrm{OOS}}$ is an expanding historical mean of realized variance computed from
  % 2026-10-22 (user, E36): parked -- Table 1 now holds the tree ensembles too. Was: "For the pooled linear models of Table~\ref{tab:rh_oos} it is the mean"
  \rhoosBenchStart{} through day $t - 1$. For the pooled models of Table~\ref{tab:rh_oos}, linear and tree
  ensembles, it is the mean
  realized variance of the same half-hour bar over every earlier trading day; for the closing-interval models of
  Table~\ref{tab:rh_oos_last30} it is the mean realized variance of the closing half hour. Both rely solely on
  information available prior to the forecast origin.
% 2026-10-15 (user): parked -- last-30-min scheme moved to the end of the economics section (this bullet is now in
% sec:rh_econ_last30, "Squared-error loss"). Was:
% \item \textbf{One day dominates the squared error of the closing-interval trees:} \rhoosBigDay{}, a day outside the
%   option-data subsample, carries \rhoosBigDaySEMin--\rhoosBigDaySEMax\% of their total squared error over
%   $\mathcal T_{\mathrm{oos}}$ but at most \rhoosBigDayQLMax\% of their total QLIKE loss, so their MSE and
%   $R^2_{\mathrm{OOS}}$ are driven by that day and their average QLIKE is not.
\end{itemize}

% 2026-10-15 (user): parked -- moved to Methods (Forecast evaluation), numbers removed
% (05_methods.tex, before Statistical tests; label sec:rh_results_eval kept there). The numbers removed and where
% they live: the out-of-sample dates and T_oos (\rhrDays) -> Experimental setup above and Data (sec:rh_data); the
% lags q_NW (\fqsDMLag, \fqsHACLag) -> Experimental setup above ("Loss and tests") and the table notes. The
% bullets "Diebold--Mariano test", "Long-run variance" and "Lag truncation" were merged into Statistical tests
% (sec:rh_methods_tests). Result statements, parked here rather than kept, as each duplicates a table: the
% benchmark forecast's average QLIKE (\rhrOlsQ; Model classes below and Table tab:rh_rank) and its sign strategy's
% annualized Sharpe ratio at mid-quotes (\fqsOlsSh; Table tab:rh_econ_control, row "OLS, HAR benchmark").
% Old subsection, verbatim:
% \subsection{Forecast evaluation}
% \label{sec:rh_results_eval}
% \begin{itemize}
% \item \textbf{Forecast and realization.} The variance forecast is the back-transformed fit,
%   $\hat\sigma^2_t = \hat y_t^{\,2}\, s^2_{t,c}$, with no bias correction. The realization is the actual realized
%   variance of the closing half hour (15:30--16:00 ET), $\RV_t \equiv \RV_{t,c}$; it is not winsorized (the winsorized
%   $\tilde y_{t,c}$ enters estimation only). The out-of-sample period is the set $\mathcal T_{\mathrm{oos}}$ of
%   $T_{\mathrm{oos}} = |\mathcal T_{\mathrm{oos}}| = \rhrDays$ trading days; in the HAC formulas, $t - j$ is the
%   out-of-sample day $j$ days before $t$.
% % 2026-10-15 (notation pass): parked -- the realization is now the actual RV (user decision #10, data audit); the old
% % bullet did not say which RV, and the scorer in the pointer below used the winsorized one.
% % \item Forecast $\hat\sigma^2_t = \hat y_t^2 B_t$ (plain back-transform); target $RV_t$, 15:30--16:00; $N = 866$
% %   trade days.
% % Code: RV_t is the winsorised target true_raw = true_adj^2 B, not the raw RV
% % (experiments/master_table_close.py:28-31, 1224-1236). The methods text says "raw RV, never winsorised".
% % 2026-10-15 (notation pass): the pointer above is the pre-switch scorer. Scoring moves to the actual RV_{t,c} (user
% % decision #10); the QLIKE, MSE and R^2 macros of this section are regenerated on it.
% \item \textbf{Loss.} The daily QLIKE loss \citep{Patton2011} and its out-of-sample average:
%   \[
%   L_t = \QLIKE(\RV_t, \hat\sigma^2_t) = \frac{\RV_t}{\hat\sigma^2_t} - \log\frac{\RV_t}{\hat\sigma^2_t} - 1,
%   \qquad
%   \overline{\QLIKE} = \bar L = \frac{1}{T_{\mathrm{oos}}}\sum_{t \in \mathcal T_{\mathrm{oos}}} L_t .
%   \]
% % 2026-10-15 (consolidation): parked -- "HAR benchmark" names the predictor set (guide section 5); the OLS on it is
% % the benchmark forecast. Old bullet: "The HAR benchmark is the closing-interval OLS regression on the HAR components
% % and the calendar dummies ... (Section~7)."
% \item \textbf{Benchmark forecast.} The closing-interval OLS regression on the HAR benchmark, the HAR components and
%   the calendar dummies \citep{Corsi2009}. Its average QLIKE is $\overline{\QLIKE} = \rhrOlsQ$, and its sign strategy,
%   which buys the 0DTE SPXW straddle at 15:30 ET when the variance forecast $\hat\sigma^2_t$ exceeds the option-implied
%   variance and sells it otherwise, has an annualized Sharpe ratio $\SR = \fqsOlsSh$ at mid-quotes
%   (Section~\ref{sec:rh_econ_strategy}).
%   % 2026-10-15 (finish): parked -- the Introduction (which defined the sign strategy) is parked, so the strategy is
%   % described at this first use. Was: "... and its sign strategy has an annualized Sharpe ratio at mid-quotes of
%   % $\SR = \fqsOlsSh$ (Section~\ref{sec:rh_econ})."
% % Code: master_table_close.py:233 sets REFERENCE = "blk2" for qlike_pct_vs_ref and dm_vs_ref; only the GUI's DM
% % (master_table_gui/rows.py:1585-1605) is against the per-bar OLS. The % and DM vs the OLS need porting.
% \item \textbf{Relative loss.} For forecast $A$ against forecast $B$ (the benchmark forecast unless stated otherwise), the loss
%   differential is $d_t = L^{A}_t - L^{B}_t$, and the percentage change in average QLIKE is
%   \[
%   \Delta\% = 100\,\frac{\bar d}{\bar L^{B}}, \qquad \bar d = \frac{1}{T_{\mathrm{oos}}}\sum_{t \in \mathcal T_{\mathrm{oos}}} d_t .
%   \]
% % 2026-10-15 (finish): parked -- the horizon symbol h dropped (h is the boosting Hessian); the forecasts are one step
% % ahead. Was: "... small-sample correction at horizon $h = 1$ \citep{HarveyLeybourneNewbold1997} ..."
% \item \textbf{Diebold--Mariano test} \citep{DieboldMariano1995}, with the Harvey--Leybourne--Newbold small-sample
%   correction for one-step-ahead forecasts \citep{HarveyLeybourneNewbold1997} and a two-sided normal $p$-value:
%   \[
%   \DM = \sqrt{\frac{T_{\mathrm{oos}}-1}{T_{\mathrm{oos}}}}\;\frac{\bar d}{\sqrt{\hat\Omega/T_{\mathrm{oos}}}},
%   \qquad
%   p = 2\bigl\{1 - \Phi(|\DM|)\bigr\}.
%   \]
% % Code: HLN factor sqrt((N + 1 - 2h + h(h-1)/N)/N) with h = 1 (src/evaluation/diebold_mariano.py:81-84).
% \item \textbf{Long-run variance:} \citet{NeweyWest1987}, with Bartlett weights:
%   \[
%   \hat\Omega = \hat\gamma_0 + 2\sum_{j=1}^{q_{\mathrm{NW}}} w_j\,\hat\gamma_j,
%   \qquad
%   w_j = 1 - \frac{j}{q_{\mathrm{NW}}+1},
%   \qquad
%   \hat\gamma_j = \frac{1}{T_{\mathrm{oos}}-j}\sum_{t=j+1}^{T_{\mathrm{oos}}}(d_t - \bar d)(d_{t-j} - \bar d).
%   \]
% % Code: dm_test (diebold_mariano.py:67) takes np.mean over the N - j pairs, so gamma_j divides by N - j; the trade's newey_west_t
% % (notebooks/atm_straddle_lib.py:1240-1261) divides by N. The two HAC conventions differ.
% % 2026-10-15 (consolidation): this display keeps the code's T_oos - j divisor; Methods (05_methods.tex) displays 1/|O|
% % pending the user's decision on FLAG (4) there. Make the two displays agree once it is decided.
% % 2026-10-15 (finish): parked -- h = 1 substituted (max(h - 1, .) is the floor term). Was:
% % $q_{\mathrm{NW}} = \max\bigl(h - 1, \lfloor 4(T_{\mathrm{oos}}/100)^{2/9}\rfloor\bigr) = \fqsDMLag$.
% \item \textbf{Lag truncation}
%   $q_{\mathrm{NW}} = \lfloor 4(T_{\mathrm{oos}}/100)^{2/9}\rfloor = \fqsDMLag$.
%   $\DM < 0$ favors forecast $A$. ``Significant'' means significant at the 5\% level ($p < 0.05$).
% \item \textbf{Squared-error and absolute-error loss,} in variance units:
%   \[
%   \mathrm{MSE} = \frac{1}{T_{\mathrm{oos}}}\sum_{t \in \mathcal T_{\mathrm{oos}}}\bigl(\RV_t - \hat\sigma^2_t\bigr)^2,
%   \qquad
%   \mathrm{MAE} = \frac{1}{T_{\mathrm{oos}}}\sum_{t \in \mathcal T_{\mathrm{oos}}}\bigl|\RV_t - \hat\sigma^2_t\bigr|.
%   \]
% \item \textbf{Out-of-sample $R^2$,} with the out-of-sample mean of the realized variance as the benchmark forecast:
%   \[
%   R^2_{\mathrm{OOS}} = 1 - \frac{\sum_{t \in \mathcal T_{\mathrm{oos}}} \bigl(\RV_t - \hat\sigma^2_t\bigr)^2}{\sum_{t \in \mathcal T_{\mathrm{oos}}} \bigl(\RV_t - \overline{\RV}\bigr)^2},
%   \qquad
%   \overline{\RV} = \frac{1}{T_{\mathrm{oos}}}\sum_{t \in \mathcal T_{\mathrm{oos}}} \RV_t .
%   \]
% % Code: master_table_gui/rows.py:1549-1563. ybar is the mean over all N scored days: look-ahead. A causal
% % benchmark would use the expanding mean of RV_{t'}, t' < t.
% \item \textbf{Calibration:} the ratio of average realized to average forecast variance, $\overline{\RV}/\overline{\hat\sigma^2}$,
%   and the frequency of under-prediction, $\frac{1}{T_{\mathrm{oos}}}\sum_{t \in \mathcal T_{\mathrm{oos}}} \mathbf{1}\{\RV_t > \hat\sigma^2_t\}$.
% \end{itemize}

\subsection{The full predictor set versus the HAR benchmark}
\label{sec:rh_results_predictors}
% 2026-10-15 (notation pass): parked -- this typed table duplicates Table~\ref{tab:rh_axis1} (generated, same numbers);
% its numbers are pre-switch (winsorized target), mix 3- and 4-decimal QLIKE, and put a p-value in the DM column, so it
% does not meet the table standard (NOTATION_GUIDE section 8). The text now points to tab:rh_axis1.
% \begin{table}[h]\centering\small
% \resizebox{\textwidth}{!}{%
% \begin{tabular}{lrrrr}
% \toprule
% Family & QLIKE, HAR + cal. & QLIKE, all & Change & DM \\
% \midrule
% Per-bar ridge & 0.0946 & 0.0958 & $+1.2\%$ & $+0.26$ \\
% Per-bar lasso & 0.0946 & 0.0914 & $-3.4\%$ & $-1.22$ \\
% Per-bar elastic net & 0.0946 & 0.0942 & $-0.4\%$ & $-0.10$ \\
% Per-bar LightGBM & 0.0986 & 0.0939 & $-4.7\%$ & $-1.43$ \\
% Per-bar XGBoost & 0.0976 & 0.0935 & $-4.2\%$ & $-1.41$ \\
% Per-bar random forest & 0.0968 & 0.0966 & $-0.2\%$ & $-0.05$ \\
% Per-bar MinGRU (QLIKE-tuned) & 0.0985 & 0.0948 & $-3.8\%$ & $p=0.34$ \\
% Pooled ridge / lasso / elastic net & 0.107 & 0.102 / 0.104 / 0.102 & $-4.3 / -1.7 / -4.7\%$ & $-1.29 / -0.50 / -1.42$ \\
% Pooled LightGBM / XGBoost / RF & 0.121--0.124 & 0.106--0.107 & $-12$ to $-14\%$ & about $-4.5$ \\
% \bottomrule
% \end{tabular}}
% \caption{All features vs HAR + calendar within each family, on 866 days (paired DM; negative favours all features).}
% \end{table}
\begin{itemize}
% 2026-10-15 (user): parked -- last-30-min scheme moved to the end of the economics section; the comparison is now
% within the pooled intraday scheme. Was: "... the same model $m$ and estimation sample $\mathcal I$ on the two
% predictor sets, ... the loss of the forecast $f_{m,\mathcal X,\mathcal I}$:"
\item \textbf{Paired comparison within each model class.} To isolate the predictive contribution of the exogenous
  % 2026-10-25 (user): parked -- the label I = {1, ..., M}; the pooled models are estimated on I_rh.
  features within each intraday model class $m$ ($\mathcal I = \mathcal I_{\mathrm{rh}}$), we form the daily loss
  differential $d_t = L^{\text{Full}}_t - L^{\text{HAR}}_t$, where
  % 2026-10-15 (user): parked -- "park all tables that I haven't mentioned in this chat". Was: $f_{m,\mathcal X,\mathcal I}$ (Table~\ref{tab:rh_axis1}):
  $L^{\mathcal X}_t$ denotes the QLIKE loss of the forecast $f_{m,\mathcal X,\{1,\dots,M\}}$. We report the
  percentage change in average loss $\Delta\%$ defined in Section~\ref{sec:rh_results_eval}, along with the
  Diebold--Mariano statistic $\DM$ computed on $d_t$ over $\mathcal T_{\mathrm{oos}}$ using lag truncation
  $q_{\mathrm{NW}} = \rhoosDMLag$.
  % 2026-10-16 (formula pass): parked -- the display repeated that relative loss and also wrote it a second way,
  % 100 * bar d / bar L^{HAR}. Was:
  % Delta% = 100 (bar L^{Full} / bar L^{HAR} - 1) = 100 * bar d / bar L^{HAR},
  % DM on d_t, t in T_oos, q_NW = \rhoosDMLag.
% Code: writeup/make_fig_qlike_sharpe.py:287-320 (dq_pct; dm_test(all, base) at the automatic lag).
% 2026-10-15 (consolidation): parked -- pre-switch numbers (winsorized target). Was: "12--14\% ($\DM$ about $-4.5$)";
% "83--91\%"; MinGRU "by 3.8\%"; "the network ties the elastic net and ridge".
% 2026-10-15 (user): parked -- on the 866 option-data days only the three pooled trees were significant; on the
% evaluation span the pooled elastic net is too (asserted in make_rehaul_results.py). Was: "... ($\DM =
% \rhrAfLgbmPoolDm$, $\rhrAfXgbPoolDm$, $\rhrAfRfPoolDm$), and only these are significant at the 5\% level."
% 2026-10-15 (user): parked -- last-30-min scheme moved to the end of the economics section; the bullet is restated
% for the pooled intraday classes (all six reductions, the linear ones added; asserted in make_rehaul_results.py).
% Was: "\textbf{The full predictor set lowers average QLIKE in most model classes} ($\Delta\% < 0$). The reductions are
% largest for the pooled intraday trees, ... ($\DM = \rhrAfLgbmPoolDm$, $\rhrAfXgbPoolDm$, $\rhrAfRfPoolDm$); these
% and the pooled intraday elastic net ($\Delta\% = \rhrAfEnetPoolPct$, $\DM = \rhrAfEnetPoolDm$) are the only
% reductions significant at the 5\% level."
% 2026-10-22 (user, E34: one design for the 12 pooled forecasts): parked -- the heading and "across every intraday model class ($\Delta\% < 0$ in all ... classes)": with the pooled linear models on the trees' design the closing-half-hour QLIKE is lower with the full set in \rhrAfPoolLower{} of the classes (make_rehaul_results.py, E34 CLAIM); the count is printed. Was: "\item \textbf{Consistent QLIKE reductions across all model classes.} Augmenting the HAR benchmark with the full predictor set reduces the average out-of-sample QLIKE loss across every intraday model class ($\Delta\% < 0$ in all \rhrAfPoolPairs{} classes), with \rhrAfPoolLowerSig{} of these reductions statistically significant at the"
  % 2026-10-22 (user, E36): parked -- the pointer: these numbers are the closing half hour over T_oos (make_rehaul_results.py, T1, parked), not Table 1, which scores every bar from 2012; Table 1's own counts added. Was: "5\% level (see Table~\ref{tab:rh_oos}):"
  % 2026-10-23 (user, final rerun, phase 2): parked -- the order and the closing-half-hour clause: on the final rerun the full set lowers the closing-half-hour QLIKE for the three tree ensembles only and raises it for the three penalized linear models; every bar first (Table 1), then the closing half hour. Was: "Augmenting the HAR benchmark with the full predictor set lowers the average out-of-sample QLIKE loss in \rhrAfPoolLower{} of the \rhrAfPoolPairs{} intraday model classes, with \rhrAfPoolLowerSig{} of these reductions statistically significant at the 5\% level (forecasts of the closing half hour over $\mathcal T_{\mathrm{oos}}$). Scored on every bar of Table~\ref{tab:rh_oos}, from \rhorSpanStart{} to \rhorSpanEnd, the full predictor set has the lower QLIKE in \rhorAfLower{} of the \rhorAfPairs{} classes, \rhorAfLowerSig{} of them significantly at the 5\% level. On the closing half hour:"
\item \textbf{QLIKE: full predictor set against the HAR benchmark.} Scored on every bar of
  Table~\ref{tab:rh_oos}, from \rhorSpanStart{} to \rhorSpanEnd, augmenting the HAR benchmark with the full
  predictor set lowers the average out-of-sample QLIKE loss in \rhorAfLower{} of the \rhorAfPairs{} intraday model
  classes, \rhorAfLowerSig{} of them significantly at the 5\% level. For the forecasts of the closing half hour over
  $\mathcal T_{\mathrm{oos}}$, it lowers the loss in \rhrAfPoolLower{} of the \rhrAfPoolPairs{} classes, the
  \rhrPoolNTreeWord{} tree ensembles, with \rhrAfPoolLowerSig{} of these reductions statistically significant at
  the 5\% level, and raises it for the \rhrPoolNLinWord{} penalized linear models:
  \begin{itemize}
  % 2026-10-22 (user, E34): parked -- with the base-5 pooled trees the tree reductions are no longer the largest or each significant (make_rehaul_results.py, E34 CLAIM); p-values and a count printed, no ranking. Was: "\item Regarding the tree ensembles, we observe the largest relative improvements: $\Delta\% = \rhrAfLgbmPoolPct$ for LightGBM ($\DM = \rhrAfLgbmPoolDm$), $\rhrAfXgbPoolPct$ for XGBoost ($\DM = \rhrAfXgbPoolDm$), and $\rhrAfRfPoolPct$ for the random forest ($\DM = \rhrAfRfPoolDm$), each of which is statistically significant at the 5\% level."
  \item Regarding the tree ensembles, the reductions are $\Delta\% = \rhrAfLgbmPoolPct$ for LightGBM
    ($\DM = \rhrAfLgbmPoolDm$, \rhrAfLgbmPoolPval), $\rhrAfXgbPoolPct$ for XGBoost ($\DM = \rhrAfXgbPoolDm$,
    \rhrAfXgbPoolPval), and $\rhrAfRfPoolPct$ for the random forest ($\DM = \rhrAfRfPoolDm$, \rhrAfRfPoolPval);
    % 2026-10-23 (phase 2): parked -- typed "three". Was: "... in \rhrAfPoolTreeLowerSig{} of the three classes."
    the reduction is statistically significant at the 5\% level in \rhrAfPoolTreeLowerSig{} of the
    \rhrPoolNTreeWord{} classes.
  % 2026-10-22 (user, E34: one design for the 12 pooled forecasts): parked -- the elastic net's "statistically significant reduction" and the ridge and lasso "reductions" (with the pooled linear models on the trees' design ridge and lasso have a higher QLIKE and no change is significant); p-values and a count printed. Was: "\item Regarding the penalized linear models, the elastic net achieves a statistically significant reduction of $\Delta\% = \rhrAfEnetPoolPct$ ($\DM = \rhrAfEnetPoolDm$), while ridge regression and the lasso yield reductions of $\rhrAfRidgePoolPct$ ($\DM = \rhrAfRidgePoolDm$, \rhrAfRidgePoolPval) and $\rhrAfLassoPoolPct$ ($\DM = \rhrAfLassoPoolDm$, \rhrAfLassoPoolPval), respectively."
  % 2026-10-23 (user, final rerun, phase 2): parked -- "the change is a statistically significant reduction ... in \rhrAfPoolLinLowerSig{} of the three classes": on the final rerun every linear change is an increase (ridge and lasso significant at 5 %, the elastic net not). Was: "\item Regarding the penalized linear models, the QLIKE changes are $\Delta\% = \rhrAfEnetPoolPct$ for the elastic net ($\DM = \rhrAfEnetPoolDm$, \rhrAfEnetPoolPval), $\rhrAfRidgePoolPct$ for ridge regression ($\DM = \rhrAfRidgePoolDm$, \rhrAfRidgePoolPval), and $\rhrAfLassoPoolPct$ for the lasso ($\DM = \rhrAfLassoPoolDm$, \rhrAfLassoPoolPval); the change is a statistically significant reduction at the 5\% level in \rhrAfPoolLinLowerSig{} of the three classes."
  \item Regarding the penalized linear models, the full predictor set raises the QLIKE loss by
    $\Delta\% = \rhrAfRidgePoolPct$ for ridge regression ($\DM = \rhrAfRidgePoolDm$, \rhrAfRidgePoolPval),
    $\rhrAfLassoPoolPct$ for the lasso ($\DM = \rhrAfLassoPoolDm$, \rhrAfLassoPoolPval), and $\rhrAfEnetPoolPct$
    for the elastic net ($\DM = \rhrAfEnetPoolDm$, \rhrAfEnetPoolPval); at the 5\% level the increase is
    statistically significant for ridge regression and the lasso and not for the elastic net.
  \end{itemize}
  % 2026-10-15 (consolidation): the \rhrAf...Pool macros (tab:rh_axis1, actual RV) replace the typed 13--15 %, -4.6.
% 2026-10-23 (user, final rerun, phase 2): added -- the robustness window (writeup/make_rehaul_robust_window.py,
% \rhrb... macros; Table tab:rh_robust_w3500 in the appendix). Checked against the per-arm macros: the significant
% window changes are the three HAR-set penalized models on every bar and the full-set ridge and lasso on every bar
% (\rhrbNSigBar), the three full-set linear models at 16:00 (\rhrbNSigClose); no tree arm changes significantly.
\item \textbf{Robustness to the window length.} With a rolling window of \rhrbWinLong{} trading days in place of
  \rhrbWinMain{} (Table~\ref{tab:rh_robust_w3500}), the full predictor set remains significantly more accurate than
  the HAR benchmark on every bar in \rhrbNFullLowerSigBarLong{} of the \rhrbNModelsLong{} model classes, and no
  pooled forecast's average QLIKE moves by more than \rhrbMaxAbsPctBar\% on every bar or \rhrbMaxAbsPctClose\% on
  the closing half hour. The changes that are statistically significant at the 5\% level (\rhrbNSigBar{} of the
  \rhrbNArms{} pooled forecasts on every bar, \rhrbNSigClose{} on the closing half hour) are all in the penalized
  linear models; for example, the closing-half-hour QLIKE of the lasso on the full predictor set changes by
  \rhrbPctLassoFullClose\% (\rhrbPLassoFullClose).
% 2026-10-15 (user): the Sharpe counts are on the option-data subsample (the straddle needs the option quotes).
% 2026-10-15 (user): parked -- last-30-min scheme moved to the end of the economics section; the model set's 51 pairs
% (mostly closing-interval tree specifications) are in sec:rh_econ_last30; the pooled counts replace them here. Was:
% \item \textbf{Across all \rhrAfPairs{} estimator pairs} of the model set, the full predictor set lowers average QLIKE
%   in \rhrAfLower{} pairs, \rhrAfLowerSig{} of them significantly (DM). On the option-data subsample
%   $\mathcal T_{\mathrm{opt}}$, it raises the annualized Sharpe ratio at mid-quotes in \fqsShHigher{} pairs,
%   \fqsShHigherSig{} of them significantly (HAC $t$-statistic on the daily return differential, Newey--West,
%   $q_{\mathrm{NW}} = \fqsHACLag$).
% 2026-10-17 (user, E16): parked -- bullet removed; the Sharpe-versus-mean concern moved to the "Full predictor set
% versus the HAR benchmark" item of sec:rh_econ_results, where the Sharpe gains are reported. Was:
% \item \textbf{Economic value of the reductions:} on the option-data subsample $\mathcal T_{\mathrm{opt}}$, the full
%   predictor set raises the annualized Sharpe ratio of the sign strategy at mid-quotes in all \rhrAfPoolPairs{}
%   classes, by \rhrAfPoolDShMin{} to \rhrAfPoolDShMax, none significantly (HAC $t$-statistic on the
%   daily return differential, Newey--West, $q_{\mathrm{NW}} = \fqsHACLag$; Section~\ref{sec:rh_econ_results}).
% \editc{orange}{E16}{Decision: move ``none significantly'' onto the mean return, or cite a Sharpe test.\\Problem: the bullet says the Sharpe increase is insignificant, and the HAC $t$ in parentheses tests the mean of the return gap $\Delta_t$. That is a different hypothesis.\\
% The bullet says the full predictor set raises the Sharpe ratio of the sign strategy in all \rhrAfPoolPairs{} classes, by \rhrAfPoolDShMin{} to \rhrAfPoolDShMax, and that none of those increases is significant. The number in parentheses is a HAC $t$-statistic on the daily return gap $\Delta_t=r^A_t-r^B_t$, with Newey--West lag $q_{\mathrm{NW}}=\fqsHACLag$. Section~\ref{sec:rh_econ_measures} builds that $t$ as a test of the mean of $\Delta_t$. The Sharpe gap $\SR(r^A)-\SR(r^B)$ divides each strategy's mean by its own standard deviation, so a test of the mean gap is a different statement from a test of the Sharpe gap. The measures section says this $t$ tests the mean. The bullet attaches the word ``significantly'' to the Sharpe gap.}
% 2026-10-15 (user): parked -- last-30-min scheme moved to the end of the economics section; the model-set MSE count
% and the closing-interval squared-error bullets are in sec:rh_econ_last30; the pooled counts replace them here. Was:
% \item \textbf{Squared-error loss disagrees:} the full predictor set raises the MSE in \rhrAfMseHigher{} of the
%   \rhrAfPairs{} pairs, \rhrAfMseHigherSig{} of them significantly.
%   \begin{itemize}
%   \item Five days in March 2020 carry \rhrSETopFiveMin--\rhrSETopFiveMax\% of
%     $\sum_{t \in \mathcal T_{\mathrm{oos}}} (\RV_t - \hat\sigma^2_t)^2$.
%   \item Excluding February--April 2020, the lasso on the full predictor set has a lower MSE than the benchmark
%     forecast.
%   \end{itemize}
% 2026-10-19 (user): parked -- "comment out the bullet on the divergence btwn qlike and squared-error loss". Was:
% \item \textbf{Divergence between QLIKE and squared-error loss.} In contrast to QLIKE, evaluation under mean
%   squared error does not confirm these reductions: incorporating the full predictor set increases the MSE in
%   \rhrAfPoolMseHigher{} of the \rhrAfPoolPairs{} model classes and lowers it in the remaining classes, with none
%   of the differences statistically significant under a Diebold--Mariano test on squared errors.
%   \begin{itemize}
%   \item This divergence reflects the extreme sensitivity of squared-error loss to a handful of volatility spikes:
%     for each pooled specification on the full predictor set, the five largest daily squared errors all occur in
%     March 2020 and account for \rhrPoolSETopFiveMin--\rhrPoolSETopFiveMax\% of the cumulative squared error
%     $\sum_{t \in \mathcal T_{\mathrm{oos}}} (\RV_t - \hat\sigma^2_t)^2$, whereas those same five days represent
%     only \rhrPoolQLTopFiveMin--\rhrPoolQLTopFiveMax\% of cumulative QLIKE loss.
%   \end{itemize}
% 2026-10-15 (user): parked -- NN content parked for now
% \item \textbf{MinGRU network} (ridge-initialized linear shortcut, trained and selected on QLIKE): the full predictor set
%   changes its average QLIKE by \rhrAfNnClosePct\% ($\DM = \rhrAfNnCloseDm$, \rhrAfNnClosePval), a reduction that is
%   not significant at the 5\% level.
%   On the full predictor set, its average QLIKE (\rhrNnAllQ) lies between those of the elastic net (\rhrEnetAllQ) and
%   ridge (\rhrRidgeAllQ).
\end{itemize}

% 2026-10-15 (user): parked -- last-30-min scheme moved to the end of the economics section. The whole subsection
% moved out: its results (pooled intraday versus closing-interval QLIKE, Table T2 macros \rhrLin/TreeBase/AllPoolPct...,
% and the Sharpe-ratio differences) are now in the final subsection of Economic value, "A final refinement: training on
% the last 30 minutes" (sec:rh_econ_last30, 07_economics.tex), which also carries the label sec:rh_results_pooled.
% Old subsection, verbatim:
% \subsection{Closing-interval versus pooled intraday models}
% \label{sec:rh_results_pooled}
% \begin{itemize}
% \item \textbf{The pooled intraday model has a higher average QLIKE than the closing-interval model in every model
%   class,} $\Delta\% = 100\bigl(\bar L^{\mathrm{pool}}/\bar L^{\mathrm{close}} - 1\bigr) > 0$, with
%   $\bar L^{\mathrm{pool}}$ and $\bar L^{\mathrm{close}}$ the average losses of $f_{m,\mathcal X,\{1,\dots,M\}}$ and
%   % 2026-10-15 (user): parked -- "park all tables that I haven't mentioned in this chat". Was: $f_{m,\mathcal X,\{c\}}$ (Table~\ref{tab:rh_axis2}):
%   $f_{m,\mathcal X,\{c\}}$:
%   \begin{itemize}
%   % 2026-10-15 (consolidation): parked -- pre-switch typed ranges; now the \rhr...PoolPct / PoolDm macros (actual RV).
%   % \item linear models: $+12$--$13\%$ on the HAR-benchmark predictors ($\DM$ about 3.4, significant at the 5\% level),
%   %   $+7$--$14\%$ on the full predictor set (significant for the lasso and the elastic net);
%   % \item trees: $+24$--$27\%$ on the HAR-benchmark predictors and $+11$--$14\%$ on the full predictor set, each
%   %   significant at the 5\% level.
%   \item linear models: $+\rhrLinBasePoolPctMin$ to $+\rhrLinBasePoolPctMax\%$ on the HAR-benchmark predictors
%     ($\DM$ from \rhrLinBasePoolDmMin{} to \rhrLinBasePoolDmMax, significant at the 5\% level),
%     % 2026-10-15 (user): parked -- on the evaluation span the elastic net's difference is not significant at 5 %
%     % (asserted in make_rehaul_results.py). Was: "(significant for the lasso and the elastic net);"
%     $+\rhrLinAllPoolPctMin$ to $+\rhrLinAllPoolPctMax\%$ on the full predictor set (significant for the lasso
%     only);
%   \item trees: $+\rhrTreeBasePoolPctMin$ to $+\rhrTreeBasePoolPctMax\%$ on the HAR-benchmark predictors and
%     $+\rhrTreeAllPoolPctMin$ to $+\rhrTreeAllPoolPctMax\%$ on the full predictor set, each significant at the 5\%
%     level.
%   \end{itemize}
% % 2026-10-15 (notation pass): the significance added to the linear (full set) and tree bullets is read from T2
% % (tab:rh_axis2) on the pre-switch scoring; re-check it after the actual-RV regeneration.
% % 2026-10-15 (user): parked -- NN content parked for now; the scope "for the linear models and the trees" existed only
% % to exclude the pooled MinGRU. Was: "\textbf{Economic value:} for the linear models and the trees, none of the
% % differences in the annualized Sharpe ratio between ..."
% \item \textbf{Economic value:} on the option-data subsample, none of the differences in the annualized Sharpe ratio
%   between the pooled intraday and the closing-interval model is significant at the 5\% level (HAC $t$-statistic on
%   the daily return differential, $q_{\mathrm{NW}} = \fqsHACLag$).
%   % 2026-10-15 (user): "on the option-data subsample" added (accuracy is on the evaluation span).
% % 2026-10-15 (notation pass): parked -- scoped to the linear models and trees: T2 (tab:rh_axis2) now carries the pooled
% % MinGRU, whose return differential on the full predictor set has HAC t = -2.30 (pre-switch), significant at 5 %.
% % \item \textbf{$\Delta\mathrm{Sharpe}$, pooled vs per-bar:} none significant.
% % 2026-10-15 (user): parked -- NN content parked for now
% % \item \textbf{Pooled intraday MinGRU network:} \todo{pending: the pooled MinGRU network, ridge-initialized and trained
% %   and selected on QLIKE}.
% % % 2026-10-15 (notation pass): T1 and T2 (tab:rh_axis1, tab:rh_axis2) already carry the pooled MinGRU (refit monthly);
% % % this bullet awaits its text.
% \end{itemize}

\subsection{Model classes}
\label{sec:rh_results_classes}
% 2026-10-15 (consolidation): parked -- pre-switch numbers and order (winsorized target). Was: "lasso 0.0914, XGBoost
% 0.0935, LightGBM 0.0939, elastic net 0.0942, MinGRU network 0.0948"; "by 3.4\% ... ($\DM = -1.22$)";
% "(average QLIKE 0.121--0.124)". On the actual RV the MinGRU ranks 5th and the benchmark forecast 9th
% (\rhrOlsRank) in Table~\ref{tab:rh_rank}.
% 2026-10-15 (user): parked -- last-30-min scheme moved to the end of the economics section. The closing-interval
% bullets below (the most accurate forecasts, the trees against the closing-interval ridge, pooled behind
% closing-interval) are now in sec:rh_econ_last30 (07_economics.tex). Old bullets, verbatim:
% \begin{itemize}
% \item \textbf{The most accurate forecasts, all closing-interval models on the full predictor set, are statistically
%   % 2026-10-15 (user): parked -- "park all tables that I haven't mentioned in this chat". Was: indistinguishable from the benchmark forecast} ($\overline{\QLIKE} = \rhrOlsQ$, rank \rhrOlsRank{} of \rhrRankN;
%   indistinguishable from the benchmark forecast} ($\overline{\QLIKE} = \rhrOlsQ$, rank \rhrOlsRank{} of \rhrRankN):
%   % 2026-10-15 (user): parked -- "park all tables that I haven't mentioned in this chat"
%   % Table~\ref{tab:rh_rank}):
%   \begin{itemize}
%   % 2026-10-15 (user): parked -- NN content parked for now; the MinGRU (5th by average QLIKE) leaves the list, so
%   % "none of the five" -> "none of the four". Was: "... LightGBM \rhrLgbmAllQ, MinGRU network \rhrNnAllQ; ... and none
%   % of the five is significantly more accurate ..."
%   % 2026-10-15 (user): parked -- the four most accurate on the 866 option-data days; on the evaluation span they are
%   % the lasso, elastic net, ridge and XGBoost on the full predictor set (asserted in make_rehaul_results.py), and
%   % LightGBM (\rhrLgbmAllQ) ranks behind the benchmark forecast. Was:
%   % \item average QLIKE: lasso \rhrLassoAllQ, XGBoost \rhrXgbAllQ, elastic net \rhrEnetAllQ, LightGBM \rhrLgbmAllQ;
%   \item average QLIKE: lasso \rhrLassoAllQ, elastic net \rhrEnetAllQ, ridge \rhrRidgeAllQ, XGBoost \rhrXgbAllQ;
%   \item relative to the benchmark forecast, the lasso changes average QLIKE by \rhrVsBmLassoCloseFullPct\%
%     ($\DM = \rhrVsBmLassoCloseFullDm$, \rhrVsBmLassoCloseFullPval), and none of the four is significantly more accurate
%     than the benchmark forecast at the 5\% level. Of the \rhrNSet{} forecasts of the model set other than the
%     benchmark forecast, \rhrNLower{} have a lower average QLIKE, \rhrNLowerSig{} of them significantly, and
%     \rhrNWorseSig{} are significantly less accurate.
%     % 2026-10-15 (consolidation): \rhrVsBmLassoCloseFull{Pct,Dm,Pval} replace the typed 3.7 % and -1.35; the count of
%     % significantly worse forecasts (\rhrNWorseSig) moved here from the table notes (results generator item 4).
%   \end{itemize}
% % 2026-10-15 (user): parked -- the campaign's counts on the 866 option-data days (\cmTrVs...); recomputed on the
% % evaluation span by make_rehaul_results.py (\rhrTrVs...; on the 866 days it reproduces 46 of 48, 21). Was:
% % "... the tree has the higher average QLIKE in \cmTrVsRidgeQAbove{} of \cmTrVsRidgeN{} comparisons,
% % \cmTrVsRidgeQAboveSig{} of them significantly."
% \item \textbf{Trees versus the closing-interval ridge on the same predictors:} the tree has the higher average QLIKE in
%   \rhrTrVsRidgeQAbove{} of \rhrTrVsRidgeN{} comparisons, \rhrTrVsRidgeQAboveSig{} of them significantly.
% \item \textbf{Pooled intraday models trail closing-interval models;} the least accurate are the pooled intraday trees on
%   the HAR-benchmark predictors (average QLIKE \rhrPoolTreeBaseQMin--\rhrPoolTreeBaseQMax).
% \end{itemize}
% 2026-10-15 (user, restructure): the pooled intraday forecasts only (the \rhrPoolN of Table tab:rh_oos, columns
% "Pooled intraday"): the ranking among them and their comparison with the benchmark forecast; every claim asserted in
% make_rehaul_results.py (the pooled-only \rhrPoolQ..., \rhrVsBmPool... macros).
\begin{itemize}
% 2026-10-18 (user, second-list 16): parked -- ranking heading; the paper is about exogenous features adding predictability of realized variance, not superiority of accuracy (bullet body is new-notes item 13). Was: "\textbf{Superior accuracy of linear specifications within each predictor set.}"
% 2026-10-22 (user, E34): parked -- with the base-5 pooled trees the pooled XGBoost on the HAR benchmark (\rhrPoolQTreeHarMin) is below the linear range (to \rhrPoolQLinHarMax), so "consistently ... on both predictor sets" no longer holds (make_rehaul_results.py, E34 CLAIM); ranges, no ranking. Table~\ref{tab:rh_oos} no longer holds the trees (note E36). Was: "Comparing across model classes in Table~\ref{tab:rh_oos}, the penalized linear models consistently achieve lower average QLIKE loss than the tree ensembles on both predictor sets:"
\item \textbf{Exogenous features and the predictability of realized variance, by model class.} Comparing across model
  classes, the average QLIKE losses of the pooled penalized linear models and of the pooled tree ensembles on
  each predictor set are:
  \begin{itemize}
  % 2026-10-22 (user, E34: one design for the 12 pooled forecasts): parked -- the ranking clause ("representing the three most accurate specifications among all \rhrPoolN{} pooled intraday forecasts"), which does not hold with the pooled linear models on the trees' design and is a ranking claim. Was: "\item Under the full predictor set, the linear models attain average QLIKE losses between \rhrPoolQLinFullMin{} and \rhrPoolQLinFullMax{}—specifically, \rhrQEnetPoolFull{} for the elastic net, \rhrQRidgePoolFull{} for ridge, and \rhrQLassoPoolFull{} for the lasso, representing the three most accurate specifications among all \rhrPoolN{} pooled intraday forecasts—compared with \rhrPoolQTreeFullMin--\rhrPoolQTreeFullMax{} for the tree ensembles."
  \item Under the full predictor set, the linear models attain average QLIKE losses between \rhrPoolQLinFullMin{}
    and \rhrPoolQLinFullMax{} (\rhrQEnetPoolFull{} for the elastic net, \rhrQRidgePoolFull{} for ridge, and
    \rhrQLassoPoolFull{} for the lasso), and the tree ensembles \rhrPoolQTreeFullMin--\rhrPoolQTreeFullMax.
  % 2026-10-22 (user, E34): parked -- "outperforming" (see the bullet above). Was: "\rhrPoolQLinHarMin--\rhrPoolQLinHarMax, outperforming the tree ensembles at"
  % 2026-10-23 (user, final rerun, phase 2): parked -- the range "\rhrPoolQLinHarMin--\rhrPoolQLinHarMax", which prints one value twice: on the HAR benchmark the three penalized models reproduce the OLS forecast at this window (frozen penalties near zero; disclosed under "Rolling re-estimation with frozen hyperparameters"). Was: "\item Under the HAR benchmark, the linear specifications achieve average QLIKE losses of \rhrPoolQLinHarMin--\rhrPoolQLinHarMax, and the tree ensembles"
  \item Under the HAR benchmark, the linear specifications achieve an average QLIKE loss of \rhrPoolQLinHarMin,
    identical for the \rhrPoolNLinWord{} of them because their HAR-benchmark penalties are numerically zero at this
    window, and the tree ensembles
    % 2026-10-17 (user, new-notes 13): parked -- no ranking claims. Was: "(which represent the least accurate pooled specifications overall)"
    \rhrPoolQTreeHarMin--\rhrPoolQTreeHarMax.
  \end{itemize}
% 2026-10-17 (user): parked -- closing-interval material (relabeled: closing-interval HAR benchmark). Was: "Comparison against the closing-interval benchmark forecast. ... than the daily closing-interval benchmark forecast"
\item \textbf{Comparison against the closing-interval HAR benchmark.} Every pooled intraday specification
  incurs a higher average QLIKE loss than the closing-interval HAR benchmark
  ($\overline{\QLIKE} = \rhrOlsQ$):
  \begin{itemize}
  \item Under the full predictor set, the average QLIKE loss of the pooled linear models exceeds that of the
    % 2026-10-17 (user): parked -- closing-interval material (relabeled: closing-interval HAR benchmark). Was: "benchmark forecast by"
    % 2026-10-22 (user, E34: one design for the 12 pooled forecasts): parked -- "although none of these differences is statistically significant at the 5\% level" (with the pooled linear models on the trees' design \rhrVsBmPoolLinFullSig{} of the three are); the count is printed.
    closing-interval HAR benchmark by $+\rhrVsBmPoolLinFullPctMin$ to $+\rhrVsBmPoolLinFullPctMax\%$, and the
    % 2026-10-22 (user, E34): parked -- with the base-5 pooled trees "each significant" and "every ... at the 1\% level" no longer hold (make_rehaul_results.py, E34 CLAIM); counts printed. Was: "differences is statistically significant at the 5\% level; for the pooled tree ensembles, the excess loss ranges from $+\rhrVsBmPoolTreeFullPctMin$ to $+\rhrVsBmPoolTreeFullPctMax\%$, each significant at the 5\% level." and "\item Under the HAR benchmark, the excess QLIKE loss of the pooled intraday models ranges from $+\rhrVsBmPoolHarPctMin$ to $+\rhrVsBmPoolHarPctMax\%$, with every difference statistically significant at the 1\% level."
    % 2026-10-23 (phase 2): parked -- typed "three" (twice). Was: "... for \rhrVsBmPoolLinFullSig{} of the three; ... for \rhrVsBmPoolTreeFullSig{} of the three."
    difference is statistically significant at the 5\% level for \rhrVsBmPoolLinFullSig{} of the
    \rhrPoolNLinWord; for the pooled tree ensembles, the excess loss
    ranges from $+\rhrVsBmPoolTreeFullPctMin$ to $+\rhrVsBmPoolTreeFullPctMax\%$, and it is significant at the 5\%
    level for \rhrVsBmPoolTreeFullSig{} of the \rhrPoolNTreeWord.
  \item Under the HAR benchmark, the excess QLIKE loss of the pooled intraday models ranges from
    $+\rhrVsBmPoolHarPctMin$ to $+\rhrVsBmPoolHarPctMax\%$; the difference is statistically significant at the
    5\% level for \rhrVsBmPoolHarSig{} of the \rhrAfPoolPairs{} models and at the 1\% level for
    \rhrVsBmPoolHarSigOne{}.
  % 2026-10-22 (user, E34: one design for the 12 pooled forecasts): parked -- "Consequently, ... narrows the accuracy gap relative to the closing-interval HAR benchmark across every intraday model class." (with the pooled linear models on the trees' design the gap narrows in \rhrVsBmPoolNarrow{} of the classes); the count is printed.
  \item Incorporating the exogenous covariates narrows the accuracy gap relative to the
    % 2026-10-17 (user): parked -- closing-interval material (relabeled: closing-interval HAR benchmark). Was: "relative to the daily closing-interval benchmark"
    % 2026-10-23 (user, final rerun, phase 2): the classes named (per-model macros \rhrVsBm<M>Pool{Har,Full}Pct: the
    % gap narrows for LightGBM, XGBoost and the forest and widens for ridge, lasso and elastic net). Was: "... in
    % \rhrVsBmPoolNarrow{} of the \rhrAfPoolPairs{} intraday model classes."
    closing-interval HAR benchmark in \rhrVsBmPoolNarrow{} of the \rhrAfPoolPairs{} intraday model classes, the
    tree ensembles, and widens it for the \rhrPoolNLinWord{} penalized linear models.
  % 2026-10-17 (user): parked -- closing-interval material (relabeled: closing-interval HAR benchmark). Was: "Because the benchmark forecast is estimated"
  \item Because the closing-interval HAR benchmark is estimated solely on the closing 30-minute interval of each day, we
    % 2026-10-17 (user): parked -- closing-interval material moved to the appendix; Section -> Appendix. Was: "extend this closing-interval training scheme to all model classes in Section~\ref{sec:rh_econ_last30}."
    extend this closing-interval training scheme to all model classes in Appendix~\ref{sec:rh_econ_last30}.
  \end{itemize}
\end{itemize}

